\documentclass[12pt,a4paper]{article}

\usepackage{fontawesome5}
\usepackage{algorithm}
\usepackage{algorithmic}
\usepackage[T1]{fontenc}
\usepackage[brazil]{babel}
\usepackage{geometry}
\usepackage{amsmath, amssymb}
\usepackage{hyperref}
\usepackage{enumitem}
\usepackage{natbib}
\usepackage{tabularx}
\usepackage{caption}
\usepackage{booktabs}
\usepackage{listings}
\usepackage{array}
\usepackage{multirow}
\usepackage[most]{tcolorbox}
\usepackage{xcolor}
\usepackage{tikz}
\usetikzlibrary{matrix, positioning, arrows.meta, shapes.geometric, trees, calc, fit, backgrounds, decorations.pathreplacing}
\usepackage{fancyhdr}
\usepackage{pgfplots}
\pgfplotsset{compat=1.18}

\setlength{\headheight}{14pt}

% ── Paleta de cores (mesma do material de AR)
\definecolor{agentClr}{RGB}{205,125,5}
\definecolor{nodeGray}{RGB}{235,234,228}
\definecolor{nodeGrayBd}{RGB}{136,135,128}
\definecolor{nodeBlue}{RGB}{230,241,251}
\definecolor{nodeBlueBd}{RGB}{24,95,165}
\definecolor{nodeRed}{RGB}{252,235,235}
\definecolor{nodeRedBd}{RGB}{163,45,45}
\definecolor{quiesBlue}{RGB}{24,95,165}
\definecolor{dangerRed}{RGB}{163,45,45}
\definecolor{evalGreen}{RGB}{40,90,10}

\renewcommand{\lstlistingname}{Listagem}
\floatname{algorithm}{Algoritmo}

\definecolor{boxblue}{RGB}{25,45,140}
\definecolor{boxbg}{RGB}{240,240,245}
\definecolor{codegray}{RGB}{248,248,248}
\definecolor{codeframe}{RGB}{180,180,200}
\definecolor{workedbg}{RGB}{232,245,232}
\definecolor{workedframe}{RGB}{40,130,40}
\definecolor{alertbg}{RGB}{255,245,220}
\definecolor{alertframe}{RGB}{180,120,0}
\definecolor{checkpointbg}{RGB}{245,232,255}
\definecolor{checkpointframe}{RGB}{100,30,160}

\definecolor{agentcolor}{RGB}{25,45,140}
\definecolor{envcolor}{RGB}{40,130,40}
\definecolor{statecolor}{RGB}{220,230,250}
\definecolor{termpos}{RGB}{180,230,180}
\definecolor{termneg}{RGB}{255,200,200}
\definecolor{blocked}{RGB}{80,80,80}
\definecolor{mlred}{RGB}{180,40,40}
\definecolor{mlblue}{RGB}{30,80,160}
\definecolor{mlgreen}{RGB}{30,120,50}
\definecolor{mlorange}{RGB}{200,100,20}
\definecolor{mlpurple}{RGB}{100,30,160}

\lstset{
  backgroundcolor=\color{codegray},
  frame=single,
  rulecolor=\color{codeframe},
  basicstyle=\ttfamily\small,
  keywordstyle=\bfseries,
  showstringspaces=false,
  breaklines=true,
  tabsize=4,
  xleftmargin=8pt,
  xrightmargin=4pt
}

\newtcolorbox[auto counter]{perguntabox}{
  colback=boxbg,
  colframe=boxblue,
  title=Pergunta~\thetcbcounter
}

\newtcolorbox{notabox}[2][]{
  enhanced,
  colback=boxbg,
  colframe=boxblue,
  boxrule=0.8pt,
  arc=2mm,
  title=\textbf{#2},
  coltitle=white,
  colbacktitle=boxblue,
  attach boxed title to top left={yshift=-2mm,xshift=2mm},
  boxed title style={boxrule=0pt, arc=1mm},
  before upper=\setlength{\parindent}{0pt},
  #1
}

\newtcolorbox{workedbox}[2][]{
  enhanced,
  colback=workedbg,
  colframe=workedframe,
  boxrule=0.8pt,
  arc=2mm,
  title=\textbf{#2},
  coltitle=white,
  colbacktitle=workedframe,
  attach boxed title to top left={yshift=-2mm,xshift=2mm},
  boxed title style={boxrule=0pt, arc=1mm},
  before upper=\setlength{\parindent}{0pt},
  #1
}

\newtcolorbox{alertabox}[2][]{
  enhanced,
  colback=alertbg,
  colframe=alertframe,
  boxrule=0.8pt,
  arc=2mm,
  title=\textbf{#2},
  coltitle=white,
  colbacktitle=alertframe,
  attach boxed title to top left={yshift=-2mm,xshift=2mm},
  boxed title style={boxrule=0pt, arc=1mm},
  before upper=\setlength{\parindent}{0pt},
  #1
}

\newtcolorbox{checkpointbox}[2][]{
  enhanced,
  colback=checkpointbg,
  colframe=checkpointframe,
  boxrule=0.6pt,
  arc=2mm,
  title=\textbf{#2},
  coltitle=white,
  colbacktitle=checkpointframe,
  attach boxed title to top left={yshift=-2mm,xshift=2mm},
  boxed title style={boxrule=0pt, arc=1mm},
  before upper=\setlength{\parindent}{0pt},
  #1
}

\geometry{margin=2.5cm}

\pagestyle{fancy}
\fancyhf{}
\fancyhead[L]{\small Bacharelado em Ciência da Computação}
\fancyhead[R]{\small GCC1734 -- Inteligência Artificial}
\fancyfoot[C]{\thepage}
\renewcommand{\headrulewidth}{0.4pt}
\renewcommand{\footrulewidth}{0pt}

\title{}
\author{}
\date{}

\begin{document}

\begin{titlepage}
\centering
\vspace*{1.8cm}
{\Large\bfseries CEFET/RJ\par}
\vspace{0.3cm}
{\large Centro Federal de Educação Tecnológica Celso Suckow da Fonseca\par}
\vspace{1.8cm}
{\large Bacharelado em Ciência da Computação\par}
\vspace{0.4cm}
{\large GCC1734 -- Inteligência Artificial\par}
\vfill
{\Huge\bfseries Aprendizado de Máquina\par}
\vspace{0.5cm}
{\LARGE Notas de Aula\par}
\vfill
{\large Prof. Eduardo Bezerra\par}
\vspace{0.2cm}
{\large Escola de Informática e Computação\par}
\vspace{0.2cm}
{\large CEFET/RJ\par}
\vfill
{\large \today\par}
\end{titlepage}

\thispagestyle{empty}
\tableofcontents
\newpage

%===========================================
\section*{Objetivos}

Ao final da leitura, você deve ser capaz de:

\medskip
\noindent\textbf{Motivação e taxonomia}
\begin{itemize}
    \item Explicar o que é aprendizado de máquina e posicioná-lo em relação a busca, planejamento e aprendizado por reforço.
    \item Distinguir os três paradigmas principais (supervisionado, não supervisionado e por reforço) e identificar o paradigma adequado a um problema dado.
    \item Enunciar a hipótese de aprendizagem e reconhecer seus limites.
\end{itemize}

\noindent\textbf{Framework de aprendizado supervisionado} $\star$
\begin{itemize}
    \item Formalizar um problema de aprendizado supervisionado como a busca de uma hipótese $h \in \mathcal{H}$ que minimiza uma função de perda $\mathcal{L}$.
    \item Distinguir erro de treinamento de erro de generalização e explicar por que minimizar o primeiro não garante o segundo.
    \item Justificar a necessidade de divisão dos dados em conjuntos de treino, validação e teste.
\end{itemize}

\noindent\textbf{Modelos lineares e gradiente descendente} $\star$
\begin{itemize}
    \item Enunciar o modelo de regressão linear e a função de perda quadrática e executar um passo de gradiente descendente manualmente.
    \item Descrever o modelo de regressão logística e interpretar sua saída como probabilidade.
    \item Explicar o papel da engenharia de características na extensão da capacidade de modelos lineares.
\end{itemize}

\noindent\textbf{Generalização, viés-variância e regularização} $\star$
\begin{itemize}
    \item Identificar sintomas de underfitting e overfitting em curvas de validação.
    \item Enunciar o dilema viés-variância e relacioná-lo à complexidade do modelo.
    \item Descrever como a regularização $\ell_1$ e $\ell_2$ penaliza modelos complexos e previne overfitting.
    \item Descrever a validação cruzada $k$-fold e explicar seu propósito.
\end{itemize}

\noindent\textbf{Árvores de decisão e métodos ensemble} $\star$
\begin{itemize}
    \item Descrever o algoritmo de construção de árvores de decisão via ganho de informação.
    \item Calcular a entropia de uma distribuição e o ganho de informação de um atributo de divisão.
    \item Explicar como Random Forests e boosting reduzem viés e variância.
\end{itemize}

\noindent\textbf{Redes neurais artificiais} $\star$
\begin{itemize}
    \item Descrever a arquitetura de um perceptron multicamadas e o papel das funções de ativação.
    \item Explicar intuitivamente como a retropropagação calcula gradientes via regra da cadeia.
    \item Conectar o treinamento de redes neurais ao gradiente descendente estocástico.
\end{itemize}

\noindent\textbf{Aprendizado não supervisionado}
\begin{itemize}
    \item Descrever o algoritmo $k$-means e executá-lo sobre um exemplo simples.
    \item Explicar intuitivamente o objetivo da análise de componentes principais (PCA).
\end{itemize}

\begin{notabox}{Roteiro de leitura}
Estas notas apresentam uma visão panorâmica do aprendizado de máquina, com foco nos conceitos e algoritmos mais centrais para a continuidade da disciplina. O material está organizado em torno de cinco aulas:

\begin{itemize}[nosep]
\item \textbf{Aula 1} (Seções 1--2): motivação, taxonomia e o framework de aprendizado supervisionado.
\item \textbf{Aula 2} (Seção 3): modelos lineares e gradiente descendente.
\item \textbf{Aula 3} (Seção 4): generalização, viés-variância, regularização e avaliação.
\item \textbf{Aula 4} (Seção 5): árvores de decisão e métodos ensemble.
\item \textbf{Aula 5} (Seções 6--7): redes neurais e aprendizado não supervisionado.
\end{itemize}

Os objetivos marcados com $\star$ são os mais centrais: concentre-se neles se o tempo for limitado. O material pressupõe familiaridade com busca, planejamento e aprendizado por reforço, abordados nas notas anteriores.
\end{notabox}

%===========================================
\section{Motivação e Contexto}
\label{sec:motivacao}

\subsection{O que é aprendizado de máquina?}

Nos módulos anteriores, estudamos agentes que agem sob diferentes graus de conhecimento sobre o ambiente. Algoritmos de busca como A* recebem um modelo completo e encontram soluções ótimas. Minimax e MCTS raciocinam sobre um modelo adversarial fornecido. O aprendizado por reforço vai além: o agente não recebe o modelo, mas aprende uma política interagindo com o ambiente e recebendo recompensas.

O \emph{aprendizado de máquina} (AM, ou \textit{machine learning}) aborda um problema diferente, embora relacionado: aprender a partir de \emph{dados}. Em vez de um modelo do ambiente ou de um ciclo de ação-recompensa, o agente recebe um conjunto de exemplos de entrada-saída (ou apenas de entradas) e deve induzir, a partir deles, uma função que generalize para novas situações.

A pergunta central do AM pode ser formulada assim: \textit{dado um conjunto de dados sobre um fenômeno, como um sistema pode extrair automaticamente regularidades que lhe permitam tomar boas decisões em situações ainda não vistas?}

Essa pergunta é relevante em uma gama surpreendentemente ampla de problemas: classificar e-mails como spam ou não spam; prever o preço de um imóvel; reconhecer a fala humana; detectar fraudes bancárias; recomendar filmes; diagnosticar doenças a partir de imagens. Em todos esses casos, seria muito difícil (frequentemente impossível) escrever manualmente um conjunto de regras que capture o fenômeno. O AM oferece uma alternativa: deixar que o algoritmo infira as regras a partir dos dados.

\subsection{Hipótese de aprendizagem}

Aprender a partir de dados só faz sentido se os dados observados carregam informação sobre casos futuros. Essa é uma suposição forte. Ela não vem do algoritmo em si; vem do problema e do modo como os dados foram coletados. O AM se apoia nessa suposição de forma explícita, na forma da \emph{hipótese de aprendizagem}.

\begin{notabox}{Hipótese de aprendizagem}
Os dados de treinamento e os dados futuros devem ser amostras do mesmo processo gerador, ou pelo menos de processos suficientemente semelhantes. Em termos práticos, espera-se que as regularidades observadas nos dados de treinamento também estejam presentes nos dados em que o modelo será usado.

Em outras palavras: \emph{o passado é informativamente semelhante ao futuro}.
\end{notabox}

Essa hipótese é o alicerce estatístico de todo o AM. Se ela for satisfeita, um modelo que aprende bem dos dados de treinamento provavelmente se sairá bem em dados novos. Se ela for violada~--- por exemplo, se os dados de treinamento forem uma amostra tendenciosa do fenômeno, ou se a distribuição mudar ao longo do tempo~--- o modelo aprendido pode generalizar mal, mesmo que seu desempenho no conjunto de treinamento seja excelente.

A hipótese de aprendizagem não é verificável em abstrato: ela depende de como os dados foram coletados e do problema em questão. Parte do trabalho de um praticante de AM é avaliar criticamente se essa hipótese é razoável para um dado conjunto de dados.

\subsection{Relação com outros paradigmas de IA}

O AM não substitui os outros paradigmas de IA; ele se combina com eles de formas precisas. A Tabela~\ref{tab:paradigmas-ia} resume as diferenças mais relevantes.

\begin{table}[ht]
\centering
\caption{Comparação entre paradigmas de IA quanto ao conhecimento disponível, ao sinal de aprendizado e ao objeto produzido.}
\label{tab:paradigmas-ia}
\renewcommand{\arraystretch}{1.3}
\begin{tabular}{p{3.0cm}p{3.2cm}p{3.0cm}p{3.8cm}}
\toprule
\textbf{Paradigma} & \textbf{Conhecimento fornecido} & \textbf{Sinal de aprendizado} & \textbf{O que é produzido} \\
\midrule
Busca (A*, UCS) & Modelo completo & Nenhum & Plano ou caminho ótimo \\
\midrule
Minimax & Modelo adversarial completo e função de utilidade & Função de utilidade & Melhor ação segundo a árvore de busca \\
\midrule
MCTS & Regras do jogo / simulador & Resultados de simulações & Estimativas de valor das ações \\
\midrule
Aprendizado por reforço & Nenhum modelo & Recompensa escalar & Política de ação \\
\midrule
\textbf{AM supervisionado} & \textbf{Nenhum modelo} & \textbf{Exemplos rotulados} & \textbf{Função entrada-saída} \\
\midrule
AM não supervisionado & Nenhum modelo & Nenhum rótulo & Estrutura nos dados \\
\bottomrule
\end{tabular}
\end{table}

Há também conexões profundas entre os paradigmas. O aprendizado por reforço com aproximação de funções (abordado nas notas anteriores) é uma instância de AM aplicada ao problema de estimar $Q^*(s,a)$: as ``características'' $f_i(s,a)$ são exatamente o tipo de representação que o AM usa, e a regra de atualização dos pesos $w_i \leftarrow w_i + \alpha \delta f_i(s,a)$ é uma instância de gradiente descendente sobre o erro temporal, exatamente o algoritmo que estudaremos na Seção~\ref{sec:lineares}.

\begin{notabox}{Por que AM agora?}
As notas de aprendizado por reforço introduziram o conceito de \emph{aproximação de funções}: em vez de uma tabela, o agente aprende um vetor de pesos que comprime o conhecimento. Aquilo era AM aplicado a AR. Nestas notas, estudamos AM de forma mais ampla: os mesmos mecanismos (representação, otimização, generalização) surgem em problemas onde o sinal de aprendizado não são recompensas, mas exemplos rotulados ou dados não rotulados.
\end{notabox}

\subsection{Taxonomia do aprendizado de máquina}

% A Figura~\ref{fig:taxonomia-am} mostra a organização dos principais paradigmas e algoritmos de AM. Esta leitura cobre os itens destacados.
A Figura~\ref{fig:taxonomia-am} situa os principais paradigmas de aprendizado de máquina no contexto da disciplina. O aprendizado supervisionado aparece como o caso em que há exemplos rotulados, isto é, pares entrada--saída usados para aprender uma função preditiva; nessa categoria entram regressão linear, regressão logística, árvores de decisão, Random Forests e redes neurais. O aprendizado não supervisionado reúne métodos que buscam estrutura nos dados sem rótulos explícitos, como agrupamentos em $k$-means e projeções em PCA. Já o aprendizado por reforço, estudado anteriormente, ocupa uma terceira categoria: nele, o agente aprende por interação com o ambiente, recebendo recompensas em vez de respostas corretas imediatas. A marca $\star$ indica os algoritmos que serão tratados nestas notas, enquanto os demais aparecem apenas para posicionamento conceitual.

\begin{figure}[ht]
\centering
\begin{tikzpicture}[
    root/.style={draw=boxblue, fill=statecolor, rounded corners=4pt, thick,
                 minimum width=4.0cm, minimum height=0.9cm,
                 align=center, font=\small\bfseries, text=boxblue},
    cat/.style={draw=envcolor, fill=workedbg, rounded corners=3pt, thick,
                minimum width=3.4cm, minimum height=0.8cm,
                align=center, font=\small\bfseries, text=envcolor},
    alg/.style={draw=black!40, fill=white, rounded corners=2pt,
                minimum width=3.0cm, minimum height=0.7cm,
                align=center, font=\scriptsize},
    algh/.style={alg, draw=boxblue!60, fill=statecolor!50, font=\scriptsize\bfseries},
    arr/.style={-{Stealth[length=5pt]}, thick, black!50}
]

\node[root] (am) at (5, 0) {Aprendizado de Máquina};

\node[cat] (sup)  at (0, -2.0) {Supervisionado};
\node[cat] (nsup) at (5, -2.0) {Não Supervisionado};
\node[cat] (ar)   at (10,-2.0) {Por Reforço};

\draw[arr] (am) -- (sup);
\draw[arr] (am) -- (nsup);
\draw[arr] (am) -- (ar);

% Supervisionado
\node[algh] at (0, -3.3) {Regressão linear $\star$};
\node[algh] at (0, -4.1) {Reg. logística $\star$};
\node[algh] at (0, -4.9) {Árvores de decisão $\star$};
\node[algh] at (0, -5.7) {Random Forests $\star$};
\node[algh] at (0, -6.5) {Redes neurais $\star$};
\node[alg]  at (0, -7.3) {SVM, $k$-NN, \ldots};

% Não supervisionado
\node[algh] at (5, -3.3) {$k$-means $\star$};
\node[algh] at (5, -4.1) {PCA $\star$};
\node[alg]  at (5, -4.9) {DBSCAN, ICA, \ldots};
\node[alg]  at (5, -5.7) {Modelos generativos};

% Por reforço
\node[alg]  at (10,-3.3) {Q-learning};
\node[alg]  at (10,-4.1) {SARSA};
\node[alg]  at (10,-4.9) {Policy gradient};
\node[alg]  at (10,-5.7) {(ver notas de AR)};

\node[font=\scriptsize, text=boxblue] at (5, -8.3)
{$\star$ = abordado nestas notas};

\end{tikzpicture}
\caption{Taxonomia parcial do aprendizado de máquina. O paradigma de aprendizado por reforço, abordado nas notas anteriores, é aqui posicionado em relação aos demais. Os itens marcados com $\star$ são cobertos nesta leitura.}
\label{fig:taxonomia-am}
\end{figure}

\subsection{Quando usar aprendizado de máquina?}

O AM é a ferramenta certa quando:

\begin{itemize}
    \item O fenômeno é complexo demais para ser descrito por regras explícitas: reconhecer faces, traduzir idiomas e entender fala são tarefas que qualquer criança executa com facilidade, mas para as quais ninguém sabe escrever um conjunto completo de regras do tipo \textit{``se a imagem tem estas bordas neste ângulo, então é um rosto''}. O espaço de variações possíveis é vasto demais, as exceções são incontáveis e as regras relevantes dependem de contexto de formas que resistem à formalização manual. Nesses casos, faz mais sentido deixar que o algoritmo infira as regularidades a partir de exemplos do que tentar enumerá-las a priori.
    \item Os dados estão disponíveis: há exemplos suficientes do comportamento desejado.
    \item O fenômeno é suficientemente estável para que a hipótese de aprendizagem seja plausível.
\end{itemize}

O AM não é a ferramenta certa quando:

\begin{itemize}
    \item O problema tem uma solução exata eficiente: um algoritmo de busca ou de programação dinâmica pode ser mais apropriado.
    \item Os dados disponíveis são insuficientes, tendenciosos ou de má qualidade.
    \item A interpretabilidade é crítica e os dados não justificam modelos opacos.
\end{itemize}

\begin{notabox}{Três perguntas antes de usar AM}
Antes de aplicar AM a um problema, pergunte:

\begin{enumerate}[nosep]
    \item Há dados suficientes e representativos do fenômeno?
    \item Existe um padrão estável que pode ser aprendido?
    \item O custo dos erros do modelo é aceitável no domínio?
\end{enumerate}

Se a resposta a qualquer uma dessas perguntas for negativa, o problema talvez precise de outra abordagem: regras explícitas, modelagem causal, simulação, coleta de dados melhor ou intervenção humana.
\end{notabox}

\begin{checkpointbox}{Micro-checkpoint}
Para cada situação abaixo, identifique o paradigma mais adequado (busca, AR, AM supervisionado, AM não supervisionado) e justifique: (a) filtrar e-mails de spam com base em exemplos rotulados; (b) ensinar um robô a andar em terrenos novos por tentativa e erro; (c) encontrar o caminho mais curto em um mapa; (d) agrupar clientes por padrão de compras, sem rótulos.
\end{checkpointbox}

\begin{perguntabox}
A hipótese de aprendizagem afirma que o passado é informativamente semelhante ao futuro. Descreva um cenário real em que essa hipótese claramente não é satisfeita e explique como isso pode comprometer um sistema de AM treinado com dados históricos. O que um praticante poderia fazer para detectar ou mitigar esse problema?
\end{perguntabox}

%===========================================
\section{O Framework de Aprendizado Supervisionado}
\label{sec:framework}

\subsection{Dados, tarefas e modelos}
\label{subsec:dados}

No aprendizado supervisionado (AS), o ponto de partida é um \emph{conjunto de dados de treinamento}:

\[
\mathcal{D} = \{(\mathbf{x}^{(1)}, y^{(1)}),\; (\mathbf{x}^{(2)}, y^{(2)}),\; \ldots,\; (\mathbf{x}^{(n)}, y^{(n)})\}
\]

onde cada $\mathbf{x}^{(i)} \in \mathbb{R}^d$ é um \emph{vetor de características} (\textit{features}) e $y^{(i)}$ é o \emph{rótulo} ou \emph{alvo} correspondente. O índice $i$ percorre os $n$ exemplos de treinamento, e $d$ é o número de características.

O objetivo do aprendizado supervisionado é encontrar uma função $h: \mathbb{R}^d \to \mathcal{Y}$, chamada \emph{hipótese}, que mapeie corretamente entradas em saídas, tanto nos dados de treinamento quanto em entradas \emph{nunca vistas antes}.

O tipo de saída $\mathcal{Y}$ define a natureza da tarefa:

\begin{itemize}
    \item \emph{Classificação}: $\mathcal{Y}$ é um conjunto finito de classes. Exemplo: $\mathcal{Y} = \{\text{spam}, \text{não-spam}\}$.
    \item \emph{Regressão}: $\mathcal{Y} = \mathbb{R}$, um valor numérico contínuo. Exemplo: preço de um imóvel em reais.
\end{itemize}

\begin{notabox}{Exemplo motivador: classificação de e-mails}
Ao longo desta seção, usaremos um problema de classificação de spam como exemplo contínuo. Cada e-mail é representado por um vetor $\mathbf{x}$ cujas componentes capturam propriedades do e-mail: frequência de certas palavras (``grátis'', ``promoção'', ``clique''), comprimento do assunto, presença de links externos, entre outras. O rótulo $y \in \{0, 1\}$ indica se o e-mail é spam ($y=1$) ou legítimo ($y=0$). O objetivo é aprender uma função $h$ que, dado um e-mail novo, preveja corretamente seu rótulo.
\end{notabox}

\subsection{O espaço de hipóteses}
\label{subsec:espaco-hipoteses}

Até aqui, falamos em aprender uma função $h$. Mas existem infinitas funções que poderiam ajustar os dados de alguma forma. Um algoritmo de AM não procura em todas elas: ele procura dentro de uma família restrita de funções. Essa família é o \emph{espaço de hipóteses}. Assim, escolher um modelo é escolher uma forma de limitar o que o algoritmo pode aprender.

Formalmente, o espaço de hipóteses $\mathcal{H}$ é um conjunto de funções parametrizadas por um vetor de parâmetros $\boldsymbol{\theta}$:

\[
\mathcal{H} = \{h_{\boldsymbol{\theta}} : \mathbb{R}^d \to \mathcal{Y} \mid \boldsymbol{\theta} \in \Theta\}
\]

A escolha de $\mathcal{H}$ é uma decisão de projeto fundamental. Ela determina:

\begin{itemize}
    \item \emph{Capacidade expressiva}: que funções $h$ podem ser representadas em $\mathcal{H}$? Uma família restrita (e.g., funções lineares) pode ser insuficiente para capturar padrões complexos. Uma família muito rica pode ``memorizar'' os dados sem generalizar.
    \item \emph{Eficiência computacional}: o problema de encontrar o melhor $h$ em $\mathcal{H}$ é tratável?
    \item \emph{Interpretabilidade}: os parâmetros $\boldsymbol{\theta}$ têm significado compreensível para humanos?
\end{itemize}

Diferentes classes de modelos correspondem a diferentes escolhas de $\mathcal{H}$: funções lineares, árvores de decisão, redes neurais, etc.

\begin{alertabox}{Não existe almoço grátis (\textit{no free lunch})}
O teorema ``no free lunch'' estabelece que, em geral, não existe um algoritmo de AM universalmente melhor que todos os outros. Todo algoritmo faz suposições implícitas sobre a estrutura do problema (o que é chamado de \textit{viés indutivo} (\textit{inductive bias}) do algoritmo). Um algoritmo eficiente para um tipo de problema pode ser ineficiente para outro.

Na prática, isso significa que a escolha do modelo (i.e., de $\mathcal{H}$) deve ser informada pelo conhecimento do domínio e validada empiricamente. Não existe o modelo ``sempre melhor''; existe o modelo mais adequado para o problema e os dados disponíveis.
\end{alertabox}

\subsection{Função de perda}
\label{subsec:perda}

Depois de escolher uma família de funções, ainda falta decidir qual hipótese dentro dessa família é a melhor. Para isso, precisamos transformar a ideia intuitiva de ``errar pouco'' em uma quantidade numérica. Essa quantidade é a \emph{função de perda} $\ell(h(\mathbf{x}), y)$, que mede o custo de prever $h(\mathbf{x})$ quando o valor verdadeiro é $y$.

Diferentes tarefas usam diferentes funções de perda:

\begin{itemize}
    \item \emph{Erro quadrático} (regressão): $\ell(\hat{y}, y) = (\hat{y} - y)^2$. Penaliza erros grandes de forma quadrática.
    \item \emph{Erro absoluto} (regressão): $\ell(\hat{y}, y) = |\hat{y} - y|$. Mais robusto a valores atípicos.
    \item \emph{Entropia cruzada binária} (classificação): $\ell(\hat{p}, y) = -y\log\hat{p} - (1-y)\log(1-\hat{p})$, onde $\hat{p}$ é a probabilidade prevista para a classe positiva. É a perda padrão para classificação logística.
    \item \emph{Perda 0-1} (classificação): $\ell(\hat{y}, y) = \mathbf{1}[\hat{y} \ne y]$. Conta simplesmente se a previsão foi errada. Simples de interpretar, mas não diferenciável. Por isso, não é adequada para otimização direta por gradiente descendente; na prática, usa-se uma perda substituta diferenciável, como a entropia cruzada.
\end{itemize}

O \emph{risco empírico} (ou custo de treinamento) é a média da perda sobre os dados de treinamento:

\[
J(\boldsymbol{\theta}) = \frac{1}{n} \sum_{i=1}^{n} \ell\!\left(h_{\boldsymbol{\theta}}(\mathbf{x}^{(i)}),\; y^{(i)}\right)
\]

O objetivo do algoritmo de treinamento é encontrar $\boldsymbol{\theta}^*$ que minimize $J(\boldsymbol{\theta})$:

\[
\boldsymbol{\theta}^* = \arg\min_{\boldsymbol{\theta}} J(\boldsymbol{\theta})
\]

\begin{notabox}{Conexão com Q-learning}
O leitor familiarizado com as notas de aprendizado por reforço reconhecerá aqui a mesma estrutura da atualização de pesos no Q-learning aproximado:
\[
w_i \leftarrow w_i + \alpha \cdot \delta \cdot f_i(s,a)
\]
onde $\delta$ é o erro temporal (diferença entre o alvo e a estimativa atual). Isso é equivalente a um passo de gradiente descendente sobre uma função perda quadrática entre o alvo de Bellman e a estimativa $Q(s,a) = \mathbf{w} \cdot \mathbf{f}(s,a)$. A única diferença é que, no AR, o alvo também depende dos pesos $\mathbf{w}$ (via $\max_{a'} Q(s',a')$), o que introduz instabilidade adicional.
\end{notabox}

\subsection{O problema fundamental: generalização}
\label{subsec:generalizacao}

A minimização do risco empírico $J(\boldsymbol{\theta})$ resolve um problema de otimização bem definido. Mas o objetivo real do AM não é desempenho nos dados de treinamento; é desempenho em dados \emph{novos}, ainda não vistos. Essa capacidade é chamada de \emph{generalização}.

Se o objetivo fosse apenas reduzir $J(\boldsymbol{\theta})$, bastaria escolher uma classe de modelos suficientemente rica e memorizar o conjunto de treino. O problema é que memorizar não é aprender. Aprender exige transferir regularidades para casos novos.

Formalmente, o que queremos minimizar é o \emph{risco esperado} (ou erro de generalização):

\[
R(\boldsymbol{\theta}) = \mathbb{E}_{(\mathbf{x}, y) \sim p_{\text{dados}}}\bigl[\ell(h_{\boldsymbol{\theta}}(\mathbf{x}), y)\bigr]
\]

onde a expectativa é tomada sobre a distribuição verdadeira dos dados $p_{\text{dados}}$. Infelizmente, $p_{\text{dados}}$ é desconhecida. Tudo o que temos é uma amostra finita $\mathcal{D}$.

% A diferença entre $R(\boldsymbol{\theta})$ (o que queremos minimizar) e $J(\boldsymbol{\theta})$ (o que realmente minimizamos) é o coração de todo o AM supervisionado. A Figura~\ref{fig:generalizacao} ilustra esse gap.
A Figura~\ref{fig:generalizacao} organiza a diferença entre otimizar um modelo e obter um modelo que generaliza. O algoritmo recebe os dados de treino $\mathcal{D}$ e ajusta os parâmetros $\boldsymbol{\theta}$ minimizando uma função de custo $J(\boldsymbol{\theta})$. Esse processo produz o modelo treinado $h_{\boldsymbol{\theta}^*}$. No entanto, o objetivo final não é apenas reduzir o erro nos exemplos usados no treinamento, mas produzir boas previsões para dados novos. Por isso, a figura separa o erro de treinamento, observado durante a otimização, do erro real, observado quando o modelo é aplicado a exemplos ainda não vistos. A diferença entre esses dois erros é o gap de generalização: quanto maior ele for, mais o modelo se ajustou aos dados de treino sem capturar padrões que se mantêm fora deles.

\begin{figure}[ht]
\centering
\begin{tikzpicture}[
    box/.style={draw, rounded corners=3pt, thick, minimum width=3.8cm,
                minimum height=1.0cm, align=center, font=\small},
    arr/.style={-{Stealth[length=6pt]}, thick},
    lbl/.style={font=\small, fill=white, inner sep=2pt}
]
\node[box, fill=statecolor, draw=agentcolor]    (train) at (0, 0) {Dados de treino\\$\mathcal{D}$};
\node[box, fill=workedbg,   draw=envcolor]      (opt)   at (5, 0) {Otimização\\$\min J(\boldsymbol{\theta})$};
\node[box, fill=boxbg,      draw=boxblue]       (model) at (10,0) {Modelo\\$h_{\boldsymbol{\theta}^*}$};
\node[box, fill=alertbg,    draw=alertframe]    (gen)   at (10,-2.5) {Dados novos\\$(\mathbf{x}_{\text{novo}}, y_{\text{novo}})$};
\node[box, fill=checkpointbg,draw=checkpointframe](gap)  at (5,-2.5)  {Gap de\\generalização};

\draw[arr, agentcolor] (train) -- (opt);
\draw[arr, envcolor]   (opt)   -- (model);
\draw[arr, boxblue]    (model) -- (gen);
\draw[arr, dashed, alertframe] (gen) -- (gap)
    node[lbl, midway, below] {erro real};
\draw[arr, dashed, envcolor]   (opt) -- (gap)
    node[lbl, midway, left]  {erro treinamento};

\node[font=\scriptsize, align=center, text=black!60] at (5,-4.0)
{O objetivo real é minimizar o erro em dados novos, não nos dados de treino.};
\end{tikzpicture}
\caption{O gap de generalização. O algoritmo de aprendizado minimiza o erro nos dados de treino, mas o que importa é o desempenho em dados novos. A diferença entre os dois é o gap de generalização, e reduzi-la é o problema central do aprendizado supervisionado.}
\label{fig:generalizacao}
\end{figure}

\begin{alertabox}{Memorizar não é generalizar}
Um modelo suficientemente flexível pode atingir erro quase zero no treino simplesmente memorizando os exemplos. Isso não significa que ele aprendeu a regularidade do fenômeno. A evidência de aprendizagem vem do desempenho em dados que não participaram do ajuste dos parâmetros.
\end{alertabox}

\subsection{Divisão dos dados: treino, validação e teste}
\label{subsec:divisao}

Como $p_{\text{dados}}$ é desconhecida, a prática padrão é estimar o erro de generalização usando dados que \emph{não foram vistos durante o treinamento}. Para isso, o conjunto de dados disponível é dividido em três partes:

\begin{enumerate}
    \item \emph{Conjunto de treinamento}: usado para ajustar os parâmetros $\boldsymbol{\theta}$ do modelo.
    \item \emph{Conjunto de validação}: usado para comparar diferentes modelos e ajustar \emph{hiperparâmetros} (parâmetros que controlam o próprio processo de aprendizado, como a complexidade do modelo ou a taxa de aprendizado).
    \item \emph{Conjunto de teste}: usado \emph{uma única vez} ao final, para estimar o desempenho real do modelo escolhido em dados novos.
\end{enumerate}

A Figura~\ref{fig:divisao-dados} resume o papel de cada subconjunto na avaliação de modelos de aprendizado de máquina. A maior parte dos dados é usada para treinamento, isto é, para ajustar os parâmetros internos do modelo, como pesos em uma regressão ou em uma rede neural. O conjunto de validação fica separado do treinamento direto, mas ainda participa do desenvolvimento: ele orienta a escolha entre modelos, hiperparâmetros e estratégias de regularização. Já o conjunto de teste deve permanecer isolado até o fim. Ele simula dados ainda não vistos e fornece uma estimativa final de desempenho depois que todas as decisões já foram tomadas. Essa separação evita que o modelo, ou o pesquisador, se adapte indiretamente aos dados usados para medir o resultado final.

\begin{figure}[ht]
\centering
\begin{tikzpicture}[
    arr/.style={-{Stealth[length=5pt]}, thick}
]
\def\total{15.5}
\def\xtrain{0}
\def\wtrain{10.15}
\def\xval{10.15}
\def\wval{2.175}
\def\xtest{12.325}
\def\wtest{3.175}

\fill[statecolor]   (\xtrain, 0) rectangle (\xtrain+\wtrain, 0.8);
\fill[workedbg]     (\xval,   0) rectangle (\xval+\wval,     0.8);
\fill[alertbg]      (\xtest,  0) rectangle (\xtest+\wtest,   0.8);

\draw[thick] (\xtrain, 0) rectangle (\total, 0.8);
\draw[thick] (\xval,   0) -- (\xval,   0.8);
\draw[thick] (\xtest,  0) -- (\xtest,  0.8);

\node[font=\small\bfseries, text=agentcolor] at (\xtrain+\wtrain/2, 0.4) {Treinamento};
\node[font=\small\bfseries, text=envcolor]   at (\xval+\wval/2,     0.4) {Val.};
\node[font=\small\bfseries, text=alertframe] at (\xtest+\wtest/2,   0.4) {Teste};

\node[font=\scriptsize, align=center, text=agentcolor] at (\xtrain+\wtrain/2, -0.6)
    {ajuste de $\boldsymbol{\theta}$};
\node[font=\scriptsize, align=center, text=envcolor] at (\xval+\wval/2, -0.6)
    {seleção do modelo};
\node[font=\scriptsize, align=center, text=alertframe] at (\xtest+\wtest/2, -0.6)
    {avaliação final};

\end{tikzpicture}
\caption{Divisão típica do conjunto de dados em treino, validação e teste. A região maior (treino) é usada para ajustar os parâmetros. A validação fica separada porque participa das escolhas de projeto: seleção de modelo e ajuste de hiperparâmetros. O teste aparece isolado porque não deve influenciar nenhuma decisão de desenvolvimento; é consultado apenas na avaliação final.}
\label{fig:divisao-dados}
\end{figure}

\begin{alertabox}{O erro mais comum em AM: vazamento de dados}
Um erro gravíssimo em AM é usar informações do conjunto de teste durante o desenvolvimento do modelo, seja para selecionar características, ajustar hiperparâmetros ou comparar alternativas. Esse erro é chamado de \emph{vazamento de dados} (\textit{data leakage}) e produz estimativas excessivamente otimistas do desempenho real.

A regra de ouro é: \emph{o conjunto de teste deve ser tratado como se não existisse até o momento da avaliação final.} Qualquer decisão de desenvolvimento deve se basear apenas nos conjuntos de treino e validação.
\end{alertabox}

\begin{checkpointbox}{Micro-checkpoint}
Explique com suas palavras a diferença entre um \emph{parâmetro} e um \emph{hiperparâmetro}. Dê um exemplo de cada. Por que o conjunto de validação é necessário, além do conjunto de teste?
\end{checkpointbox}

\begin{perguntabox}
Um estudante treina um classificador de spam com 10.000 e-mails. Ele obtém 98\% de acurácia nos dados de treinamento. Quando coloca o modelo em produção, o desempenho cai para 72\%. Identifique pelo menos três causas possíveis para essa discrepância. Para cada causa, proponha uma forma de detectá-la antes da implantação.
\end{perguntabox}

%===========================================
\section{Modelos Lineares}
\label{sec:lineares}

\subsection{Regressão linear}
\label{subsec:reg-linear}

O modelo mais simples de AM supervisionado é a \emph{regressão linear}. Dado um vetor de características $\mathbf{x} = [x_1, x_2, \ldots, x_d]^\top$, o modelo prevê:

\[
h_{\mathbf{w},b}(\mathbf{x}) = w_1 x_1 + w_2 x_2 + \cdots + w_d x_d + b = \mathbf{w}^\top \mathbf{x} + b
\]

onde $\mathbf{w} = [w_1, \ldots, w_d]^\top$ é o vetor de pesos e $b$ é o \emph{viés} (\textit{bias}), um parâmetro escalar que desloca a previsão.

\begin{notabox}{Convenção: absorvendo o viés}
É conveniente absorver $b$ no vetor de pesos adicionando uma componente constante $x_0 = 1$ ao vetor de características: $\mathbf{x} \leftarrow [1, x_1, x_2, \ldots, x_d]^\top$ e $\mathbf{w} \leftarrow [b, w_1, \ldots, w_d]^\top$. Com isso, $h_\mathbf{w}(\mathbf{x}) = \mathbf{w}^\top \mathbf{x}$, e toda análise é feita em termos de um único vetor de parâmetros $\mathbf{w} \in \mathbb{R}^{d+1}$. Essa convenção é padrão na literatura e simplifica a notação.
\end{notabox}

Para regressão linear, a função de perda padrão é o \emph{erro quadrático médio} (\textit{mean squared error}, MSE):

\[
J(\mathbf{w}) = \frac{1}{n} \sum_{i=1}^{n} \left(h_\mathbf{w}(\mathbf{x}^{(i)}) - y^{(i)}\right)^2 = \frac{1}{n} \sum_{i=1}^n \left(\mathbf{w}^\top \mathbf{x}^{(i)} - y^{(i)}\right)^2
\]

Existe uma solução analítica fechada para o vetor ótimo de pesos, chamada \emph{equação normal}:

\[
\mathbf{w}^* = (X^\top X)^{-1} X^\top \mathbf{y}
\]

onde $X$ é a matriz de dados ($n \times (d+1)$, uma linha por exemplo) e $\mathbf{y}$ é o vetor de rótulos.

\begin{alertabox}{Quando a equação normal falha}
A equação normal requer inverter a matriz $X^\top X$, de tamanho $(d+1) \times (d+1)$. Isso tem custo $O(d^3)$, o que a torna impraticável para $d$ grande (e.g., $d = 10^6$ em problemas de processamento de linguagem natural). Além disso, $X^\top X$ pode ser singular (não invertível) se as características forem linearmente dependentes. O gradiente descendente, apresentado a seguir, contorna ambos os problemas.
\end{alertabox}

\subsubsection{Exemplo resolvido}

Suponha um problema de regressão com apenas uma característica ($d=1$): prever o preço de um imóvel ($y$, em mil reais) em função de sua área ($x$, em m²). Dados de treinamento:

\medskip
\begin{center}
\begin{tabular}{ccc}
\toprule
Exemplo & $x$ (área, m²) & $y$ (preço, R\$ mil) \\
\midrule
1 & 50  & 200 \\
2 & 100 & 350 \\
3 & 150 & 500 \\
\bottomrule
\end{tabular}
\end{center}

\medskip
Com $\mathbf{w} = [b, w_1]^\top$ e $x_0 = 1$, temos $h(x) = b + w_1 x$.

\paragraph{Passo 1: verificar um candidato.} Suponha $b = 50$ e $w_1 = 3$, ou seja, $h(x) = 50 + 3x$.

\begin{align*}
h(50)  &= 50 + 3 \times 50 = 200 \quad (\text{erro} = 0) \\
h(100) &= 50 + 3 \times 100 = 350 \quad (\text{erro} = 0) \\
h(150) &= 50 + 3 \times 150 = 500 \quad (\text{erro} = 0)
\end{align*}

Neste caso, os dados são exatamente lineares, então $J(\mathbf{w}) = 0$ é atingível.

\paragraph{Passo 2: verificar um candidato ruim.} Suponha $b = 0$ e $w_1 = 4$:

\begin{align*}
J(\mathbf{w}) &= \frac{1}{3}\bigl[(200-200)^2 + (400-350)^2 + (600-500)^2\bigr] \\
               &= \frac{1}{3}\bigl[0 + 2500 + 10000\bigr] = 4166{,}7
\end{align*}

O MSE positivo indica que este não é o modelo ótimo.

\subsection{Gradiente descendente}
\label{subsec:gd}

Para minimizar $J(\mathbf{w})$, o algoritmo mais geral é o \emph{gradiente descendente}. A ideia é simples: a partir de um ponto inicial $\mathbf{w}^{(0)}$, repetidamente mova-se na direção oposta ao gradiente de $J$, que é a direção de maior descida:

\begin{center}
\begin{tcolorbox}[colback=boxbg, colframe=boxblue, width=0.92\textwidth]
\[
\mathbf{w} \;\leftarrow\; \mathbf{w} - \alpha \,\nabla_\mathbf{w} J(\mathbf{w})
\]
onde $\alpha > 0$ é a \emph{taxa de aprendizado} e $\nabla_\mathbf{w} J(\mathbf{w})$ é o gradiente de $J$ em relação a $\mathbf{w}$.
\end{tcolorbox}
\end{center}

Para regressão linear com MSE, o gradiente em relação ao peso $w_j$ é:

\[
\frac{\partial J}{\partial w_j} = \frac{2}{n} \sum_{i=1}^n \left(\mathbf{w}^\top \mathbf{x}^{(i)} - y^{(i)}\right) x_j^{(i)}
\]

A regra de atualização para cada peso $w_j$ é, portanto:

\[
w_j \;\leftarrow\; w_j - \alpha \cdot \frac{2}{n} \sum_{i=1}^n \underbrace{\left(\mathbf{w}^\top \mathbf{x}^{(i)} - y^{(i)}\right)}_{\text{erro na previsão}} x_j^{(i)}
\]

Isso tem uma interpretação intuitiva: se o modelo está superestimando ($\hat{y} > y$) e a característica $x_j$ foi positiva, então $w_j$ foi muito alto e deve ser reduzido.

\begin{notabox}{Escala das características}
O gradiente descendente é sensível à escala das características. Se uma variável varia entre 0 e 10.000 e outra entre 0 e 1, os pesos correspondentes podem ser atualizados em ritmos muito diferentes, tornando a convergência lenta ou instável. Por isso, é comum padronizar ou normalizar as características antes do treinamento~--- por exemplo, subtraindo a média e dividindo pelo desvio padrão.
\end{notabox}

\begin{notabox}{Analogia com Q-learning}
Comparando a regra de atualização acima com a do Q-learning aproximado:
\[
w_j \leftarrow w_j + \alpha \cdot \underbrace{\delta}_{\text{erro TD}} \cdot f_j(s,a)
\]
a estrutura é idêntica: ajuste proporcional ao erro (TD ou de regressão) vezes o valor da característica. No Q-learning, $\delta = [r + \gamma \max_{a'} Q(s',a')] - Q(s,a)$; na regressão linear, $\delta = y^{(i)} - \hat{y}^{(i)}$. Ambos minimizam um erro quadrático entre uma previsão e um alvo.
\end{notabox}

\paragraph{Variantes de gradiente descendente.}

Na prática, calcular o gradiente sobre \emph{todos} os $n$ exemplos a cada passo pode ser custoso. Há três variantes:

\begin{itemize}
    \item \emph{Gradiente descendente em lote (\textit{batch gradient descent})}: usa todos os $n$ exemplos. Atualização precisa, mas custosa para $n$ grande.
    \item \emph{Gradiente descendente estocástico (\textit{stochastic gradient descent}, SGD)}: usa um único exemplo por passo, escolhido aleatoriamente. Atualização ruidosa, mas muito rápida. Convergência instável, mas frequentemente suficiente.
    \item \emph{Gradiente descendente em mini-lote (\textit{mini-batch gradient descent})}: usa um subconjunto de $b$ exemplos por passo ($b$ típico: 32--256). Compromisso prático. É o padrão em aprendizado profundo.
\end{itemize}

A Figura~\ref{fig:gd} ilustra a dinâmica do gradiente descendente em lote no espaço dos parâmetros $(w_1,w_2)$. As elipses representam curvas de nível da função de custo $J(\mathbf{w})$: pontos sobre a mesma curva têm o mesmo valor de erro. O ponto $\mathbf{w}^{(0)}$ indica a inicialização dos parâmetros, e as setas mostram as atualizações sucessivas feitas pelo algoritmo. Em cada passo, o gradiente descendente usa todos os exemplos de treinamento para calcular a direção média de maior crescimento da função de custo e, em seguida, move os parâmetros na direção oposta, isto é, na direção de maior descida. À medida que as atualizações avançam, a trajetória se aproxima do mínimo $\mathbf{w}^*$, onde o custo é menor e novas atualizações passam a ter magnitude cada vez menor.

\begin{figure}[ht]
\centering
\begin{tikzpicture}[
    arr/.style={-{Stealth[length=5pt]}, thick, blue!60}
]
\begin{axis}[
    width=10cm, height=7cm,
    xlabel={$w_1$}, ylabel={$w_2$},
    xlabel style={font=\small},
    ylabel style={font=\small},
    xtick=\empty, ytick=\empty,
    xmin=-3, xmax=3, ymin=-3, ymax=3,
    axis lines=center,
    clip=false
]
% Contornos da função de custo (elipses)
\foreach \r in {0.4, 0.9, 1.5, 2.2} {
    \addplot[domain=0:360, samples=60, gray!50, line width=0.5pt]
    ({1.5*\r*cos(x)}, {0.6*\r*sin(x)});
}
% Mínimo
\node[circle, fill=envcolor, minimum size=5pt, inner sep=0] at (axis cs: 0, 0) {};
\node[font=\scriptsize, text=envcolor, right=3pt] at (axis cs: 0, 0) {$\mathbf{w}^*$};
% Trajetória do GD
\draw[arr, thick, agentcolor] (axis cs: 2.4, 1.0) -- (axis cs: 1.4, 0.6);
\draw[arr, thick, agentcolor] (axis cs: 1.4, 0.6) -- (axis cs: 0.7, 0.3);
\draw[arr, thick, agentcolor] (axis cs: 0.7, 0.3) -- (axis cs: 0.25, 0.1);
\draw[arr, thick, agentcolor] (axis cs: 0.25, 0.1) -- (axis cs: 0.05, 0.02);
\node[font=\scriptsize, text=agentcolor] at (axis cs: 2.4, 1.3) {$\mathbf{w}^{(0)}$};
\node[font=\scriptsize, text=agentcolor, align=left] at (axis cs: -2, -2)
{trajetória do GD em lote};
\end{axis}
\end{tikzpicture}
\caption{Trajetória do gradiente descendente em lote num espaço bidimensional de parâmetros. As curvas de nível representam o custo $J(\mathbf{w})$; a trajetória parte de $\mathbf{w}^{(0)}$ e converge para o mínimo $\mathbf{w}^*$, movendo-se sempre na direção de maior descida.}
\label{fig:gd}
\end{figure}

\subsubsection{Exemplo resolvido: um passo de gradiente descendente}

Continuando o exemplo anterior com $d=1$: $h(x) = b + w_1 x$. Suponha que os pesos atuais sejam $b = 0$, $w_1 = 4$, e $\alpha = 0{,}0001$.

\paragraph{Passo 1: calcular os erros.}
\begin{align*}
e_1 &= h(50)  - 200 = 200 - 200 = 0 \\
e_2 &= h(100) - 350 = 400 - 350 = +50 \\
e_3 &= h(150) - 500 = 600 - 500 = +100
\end{align*}

\paragraph{Passo 2: calcular os gradientes.}
\[
\frac{\partial J}{\partial w_1} = \frac{2}{3}(0 \cdot 50 + 50 \cdot 100 + 100 \cdot 150) = \frac{2}{3}(0 + 5000 + 15000) = \frac{40000}{3} \approx 13333
\]
\[
\frac{\partial J}{\partial b} = \frac{2}{3}(0 + 50 + 100) = 100
\]

\paragraph{Passo 3: atualizar os pesos.}
\[
w_1 \leftarrow 4 - 0{,}0001 \times 13333 \approx 4 - 1{,}33 = 2{,}67
\]
\[
b \leftarrow 0 - 0{,}0001 \times 100 = -0{,}01
\]

\paragraph{Interpretação:} $w_1$ desceu de $4$ para $2{,}67$ (em direção ao valor ótimo $3$), pois o modelo estava superestimando os preços de imóveis maiores. O processo se repete até convergência.

\subsection{Regressão logística e classificação binária}
\label{subsec:logistica}

Para classificação binária ($y \in \{0, 1\}$), um modelo linear simples $h(\mathbf{x}) = \mathbf{w}^\top \mathbf{x}$ não é adequado como probabilidade: ele pode produzir valores negativos ou maiores que $1$. A \emph{regressão logística} aplica a função \emph{sigmoide} $\sigma(z) = \frac{1}{1 + e^{-z}}$ à saída linear para forçar o resultado ao intervalo $(0, 1)$:

\[
h_\mathbf{w}(\mathbf{x}) = \sigma(\mathbf{w}^\top \mathbf{x}) = \frac{1}{1 + e^{-\mathbf{w}^\top \mathbf{x}}}
\]

O valor $h_\mathbf{w}(\mathbf{x}) \in (0,1)$ é interpretado como a probabilidade de $y = 1$ dado $\mathbf{x}$. A decisão final é obtida por limiarização:

\[
\hat{y} = \begin{cases} 1 & \text{se } h_\mathbf{w}(\mathbf{x}) \geq 0{,}5 \\ 0 & \text{caso contrário} \end{cases}
\]

O limiar $0{,}5$ corresponde a $\mathbf{w}^\top \mathbf{x} = 0$, o \emph{hiperplano de decisão} que separa as duas classes.

A Figura~\ref{fig:sigmoide} mostra como a regressão logística transforma o escore linear $z = \mathbf{w}^\top \mathbf{x}$ em uma quantidade interpretável como probabilidade. Valores muito negativos de $z$ são comprimidos para perto de $0$, favorecendo a classe negativa; valores muito positivos são comprimidos para perto de $1$, favorecendo a classe positiva. A transição ocorre de forma suave em torno de $z=0$, ponto em que $\sigma(z)=0{,}5$. Por isso, quando se adota o limiar padrão de decisão, exemplos com $z \leq 0$ são classificados como $y=0$, enquanto exemplos com $z>0$ são classificados como $y=1$. Assim, a fronteira de decisão da regressão logística continua sendo linear no espaço de entrada, mas sua saída passa a ter interpretação probabilística.

\begin{figure}[ht]
\centering
\begin{tikzpicture}
\begin{axis}[
    width=9cm, height=5cm,
    xlabel={$z = \mathbf{w}^\top \mathbf{x}$},
    ylabel={$\sigma(z)$},
    xlabel style={font=\small},
    ylabel style={font=\small},
    xmin=-6, xmax=6,
    ymin=-0.05, ymax=1.1,
    xtick={-6,-4,-2,0,2,4,6},
    ytick={0,0.5,1},
    grid=both,
    grid style={gray!20},
    clip=false
]
\addplot[thick, agentcolor, domain=-6:6, samples=100] {1/(1+exp(-x))};
\addplot[dashed, gray, domain=-6:6] {0.5};
\node[font=\scriptsize, text=gray] at (axis cs: 5, 0.55) {$0{,}5$};
% CORRIGIDO: rótulos nas posições corretas (y=0 quando z<0, sigma~0; y=1 quando z>0, sigma~1)
\node[font=\scriptsize, text=mlred,  align=left] at (axis cs: -4, 0.15) {classe $y=0$\\($z \leq 0$)};
\node[font=\scriptsize, text=envcolor, align=left] at (axis cs: 3.25, 0.8) {classe $y=1$\\($z > 0$)};
\node[circle, fill=agentcolor, minimum size=4pt, inner sep=0] at (axis cs: 0, 0.5) {};
\end{axis}
\end{tikzpicture}
\caption{A função sigmoide $\sigma(z)$, usada na regressão logística. Ela comprime qualquer valor real $z = \mathbf{w}^\top \mathbf{x}$ para o intervalo $(0,1)$, produzindo uma estimativa de probabilidade. Para $z \ll 0$, a saída se aproxima de 0 (classe negativa); para $z \gg 0$, se aproxima de 1 (classe positiva). O limiar de decisão padrão é $0{,}5$, que corresponde a $z = 0$.}
\label{fig:sigmoide}
\end{figure}

A função de perda adequada para regressão logística é a \emph{entropia cruzada binária}:

\[
J(\mathbf{w}) = -\frac{1}{n} \sum_{i=1}^n \left[ y^{(i)} \log h(\mathbf{x}^{(i)}) + (1 - y^{(i)}) \log(1 - h(\mathbf{x}^{(i)})) \right]
\]

Essa função penaliza fortemente previsões confiantes e erradas. Se o modelo prevê $h(\mathbf{x}) = 0{,}99$ (muito confiante em $y=1$) mas o rótulo verdadeiro é $y=0$, a perda é $-\log(0{,}01) \approx 4{,}6$, que é muito alta. O gradiente descendente pode ser aplicado diretamente para minimizar esta perda.

\begin{notabox}{Regressão logística não é regressão no sentido usual}
O nome ``regressão logística'' é historicamente enganoso: trata-se de um modelo de \emph{classificação}, não de regressão. A confusão ocorre porque o modelo estima uma probabilidade contínua (regressão) e depois limiariza para produzir uma classificação binária. A denominação persiste por razões históricas.
\end{notabox}

\subsection{Engenharia de características}
\label{subsec:feat-eng}

Modelos lineares têm limitações óbvias: eles só conseguem representar fronteiras de decisão lineares (hiperplanos). No entanto, a maioria dos problemas reais envolve relações não lineares. Como lidar com isso mantendo a tratabilidade dos modelos lineares?

A resposta é a \emph{engenharia de características}: transformar as variáveis originais em um novo espaço de representação onde as relações sejam (mais próximas de) lineares. Algumas transformações comuns:

\begin{itemize}
    \item \emph{Características polinomiais}: adicionar $x_1^2$, $x_1 x_2$, $x_2^2$, etc. Permite capturar relações polinomiais.
    \item \emph{Logaritmo}: $\log x_j$ para variáveis com distribuições muito assimétricas (e.g., rendas, preços).
    \item \emph{Interações}: $x_j \cdot x_k$ captura relações de multiplicação entre características.
    \item \emph{Binarização}: converter variáveis categóricas em variáveis binárias (\textit{one-hot encoding}).
\end{itemize}

\begin{notabox}{Engenharia de características e Q-learning}
Os leitores das notas de AR já encontraram engenharia de características no Q-learning com aproximação linear: as funções $f_i(s,a)$ são exatamente características projetadas manualmente para capturar aspectos relevantes do par estado-ação. O princípio é o mesmo: um modelo linear no espaço de características pode representar relações não lineares no espaço original.
\end{notabox}

\begin{alertabox}{O risco de características polinomiais: overfitting}
Adicionar características polinomiais de grau alto permite que o modelo se ajuste perfeitamente aos dados de treinamento, interpolando por todos os pontos. Mas isso geralmente prejudica a generalização, pois o modelo ``memoriza'' os dados em vez de aprender a tendência subjacente. Esse fenômeno é chamado de \emph{overfitting} e é estudado em detalhes na próxima seção.
\end{alertabox}

\begin{checkpointbox}{Micro-checkpoint}
Explique por que a regressão linear pura não consegue classificar pontos que não são linearmente separáveis. Descreva uma transformação de características que permitiria a um modelo linear separar os pontos $\{(1,1), (-1,-1)\}$ (classe $+$) dos pontos $\{(1,-1), (-1,1)\}$ (classe $-$). (\textit{Dica}: observe que esses pontos são separáveis por uma fronteira do tipo $x_1 \cdot x_2 = 0$.)
\end{checkpointbox}

\begin{perguntabox}
Um modelo de regressão linear treinado para prever o consumo de energia elétrica de edifícios usa como características a área útil e o número de andares. Após treinamento, o modelo obtém MSE = 250 no conjunto de treino e MSE = 1800 no conjunto de teste. Identifique o problema e proponha três ações concretas (relacionadas aos dados, ao modelo ou ao processo de avaliação) para melhorar o desempenho de generalização.
\end{perguntabox}

%===========================================
\section{Generalização: Viés, Variância e Regularização}
\label{sec:generalizacao}

\subsection{Underfitting e overfitting}
\label{subsec:fitting}

Toda a dificuldade do AM supervisionado pode ser resumida em um único dilema: um modelo muito simples não captura os padrões dos dados (\emph{underfitting}); um modelo muito complexo captura até o ruído dos dados de treinamento e não generaliza (\emph{overfitting}).

A Figura~\ref{fig:fitting} compara três comportamentos típicos de modelos de regressão diante do mesmo conjunto de dados. No painel à esquerda, a reta horizontal representa um modelo simples demais: ele ignora a tendência crescente dos pontos e comete erros sistemáticos, caracterizando \textit{underfitting} ou alto viés. No painel central, a reta inclinada captura a estrutura principal dos dados sem tentar explicar cada pequena flutuação individual; esse é o caso desejável, no qual há um equilíbrio entre simplicidade e capacidade de ajuste. No painel à direita, a curva passa exatamente pelos exemplos observados, mas faz isso seguindo também as irregularidades locais dos dados. Esse comportamento caracteriza \textit{overfitting}: o erro no conjunto de treino pode ser muito baixo, mas a previsão tende a ser instável em novos exemplos, pois o modelo aprendeu ruído junto com o padrão.

\begin{figure}[ht]
\centering
\resizebox{\textwidth}{!}{%
\begin{tikzpicture}
\def\ptset{
    (0.5, 1.8), (1.0, 2.3), (1.5, 2.9), (2.0, 3.2),
    (2.5, 3.8), (3.0, 4.6), (3.5, 4.0), (4.0, 5.1),
    (4.5, 5.5), (5.0, 5.3)
}

% Underfitting
\begin{scope}[xshift=0cm]
\draw[->] (0,0) -- (5.5,0) node[right, font=\tiny] {$x$};
\draw[->] (0,0) -- (0,6.5) node[above, font=\tiny] {$y$};
\foreach \p in \ptset { \fill[agentcolor] \p circle[radius=2pt]; }
% Linha plana (underfitting)
\draw[thick, mlred, domain=0.2:5.2] plot (\x, {3.8});
\node[font=\scriptsize\bfseries, text=mlred] at (2.75, 6.2) {Underfitting};
\node[font=\tiny, text=black!60, align=center] at (2.75, -0.5) {modelo simples demais\\(alto viés)};
\end{scope}

% Ajuste adequado
\begin{scope}[xshift=6.5cm]
\draw[->] (0,0) -- (5.5,0) node[right, font=\tiny] {$x$};
\draw[->] (0,0) -- (0,6.5) node[above, font=\tiny] {$y$};
\foreach \p in \ptset { \fill[agentcolor] \p circle[radius=2pt]; }
\draw[thick, envcolor, domain=0.2:5.2] plot (\x, {0.8*\x + 1.4});
\node[font=\scriptsize\bfseries, text=envcolor] at (2.75, 6.2) {Ajuste adequado};
\node[font=\tiny, text=black!60, align=center] at (2.75, -0.5) {equilíbrio entre\\viés e variância};
\end{scope}

% Overfitting
\begin{scope}[xshift=13cm]
\draw[->] (0,0) -- (5.5,0) node[right, font=\tiny] {$x$};
\draw[->] (0,0) -- (0,6.5) node[above, font=\tiny] {$y$};
\foreach \p in \ptset { \fill[agentcolor] \p circle[radius=2pt]; }
% Curva que passa exatamente pelos pontos (overfitting)
\draw[thick, mlred, domain=0.4:5.1]
    plot[smooth] coordinates {
        (0.5, 1.8) (1.0, 2.3) (1.5, 2.9) (2.0, 3.2)
        (2.5, 3.8) (3.0, 4.6) (3.5, 4.0) (4.0, 5.1)
        (4.5, 5.5) (5.0, 5.3)
    };
\node[font=\scriptsize\bfseries, text=mlred] at (2.75, 6.2) {Overfitting};
\node[font=\tiny, text=black!60, align=center] at (2.75, -0.5) {modelo complexo demais\\(alta variância)};
\end{scope}

\end{tikzpicture}}
\caption{Três casos de ajuste em regressão. À esquerda, o modelo é simples demais (linha plana) e não captura a tendência crescente dos dados: underfitting. No centro, o modelo é linear e captura a tendência sem se prender ao ruído: ajuste adequado. À direita, o modelo passa exatamente por todos os pontos mas oscilará muito em dados novos: overfitting.}
\label{fig:fitting}
\end{figure}

% A Figura~\ref{fig:curvas-complexidade} mostra as curvas de erro típicas de treino e validação em função da complexidade do modelo.
A Figura~\ref{fig:curvas-complexidade} mostra como o erro varia à medida que aumentamos a complexidade do modelo. Modelos muito simples, à esquerda, não conseguem capturar os padrões relevantes dos dados: tanto o erro de treino quanto o erro de validação permanecem altos, caracterizando \textit{underfitting}. À medida que a complexidade aumenta, o modelo passa a representar melhor a relação entre entradas e saídas, e o erro de validação diminui. A região central representa o melhor compromisso: o modelo é expressivo o bastante para aprender padrões úteis, mas ainda não está ajustado demais às particularidades do conjunto de treino. À direita, modelos muito complexos continuam reduzindo o erro de treino, pois conseguem se adaptar cada vez mais aos exemplos observados. No entanto, o erro de validação volta a crescer, indicando \textit{overfitting}: o modelo deixa de aprender apenas padrões gerais e passa a capturar ruídos ou particularidades do treino. Por isso, a escolha da complexidade deve ser guiada pelo erro de validação, não pelo erro de treino.

\begin{figure}[ht]
\centering
\begin{tikzpicture}
\begin{axis}[
    width=10cm, height=6cm,
    xlabel={Complexidade do modelo},
    ylabel={Erro},
    xlabel style={font=\small},
    ylabel style={font=\small},
    xmin=0, xmax=10,
    ymin=0, ymax=3,
    xtick=\empty, ytick=\empty,
    clip=false,
    legend pos=north east,
    legend style={font=\scriptsize}
]
% Zonas (desenhadas antes das curvas para ficarem atrás)
\fill[mlred!10]   (axis cs: 0,   0) rectangle (axis cs: 2.5, 3);
\fill[envcolor!10](axis cs: 2.5, 0) rectangle (axis cs: 5.5, 3);
\fill[mlred!10]   (axis cs: 5.5, 0) rectangle (axis cs: 10,  3);

% Treino: decresce monotonicamente em direção a ~0
\addplot[thick, agentcolor, domain=0.5:9.5, samples=100]
    {0.03 + 1.5*exp(-0.52*x)};
\addlegendentry{erro de treino}

% Validação: formato de U, sempre acima do treino
\addplot[thick, mlred, domain=0.5:9.5, samples=100]
    {0.45 + 1.05*exp(-0.42*x) + 0.022*(x-3)^2};
\addlegendentry{erro de validação}

% Rótulos das zonas posicionados abaixo das curvas
\node[font=\scriptsize, text=mlred] at (axis cs: 1.25, 1.6) {underfitting};
\node[font=\scriptsize, text=envcolor] at (axis cs: 4.0,  1.6) {zona adequada};
\node[font=\scriptsize, text=mlred] at (axis cs: 7.75, 1.6) {overfitting};
\end{axis}
\end{tikzpicture}
\caption{Curvas de erro em função da complexidade do modelo. O erro de treino decresce monotonicamente e se aproxima de zero (o modelo memoriza os dados). O erro de validação tem formato de~U: decresce até um ponto ótimo, depois sobe (overfitting). A complexidade ótima é aquela que minimiza o erro de validação.}
\label{fig:curvas-complexidade}
\end{figure}

\subsection{O dilema viés-variância}
\label{subsec:bias-variance}

O dilema viés-variância é a formalização matemática do problema de underfitting e overfitting. Para qualquer modelo e qualquer ponto de teste $\mathbf{x}$, o erro quadrático esperado pode ser decomposto como:

\[
\mathbb{E}\left[\left(h(\mathbf{x}) - y\right)^2\right] = \underbrace{\text{Viés}^2(h(\mathbf{x}))}_{\text{erro sistemático}} + \underbrace{\text{Variância}(h(\mathbf{x}))}_{\text{sensibilidade ao dado}} + \underbrace{\sigma^2_{\text{ruído}}}_{\text{irreducível}}
\]

onde a expectativa é tomada sobre diferentes conjuntos de treinamento amostrados da mesma distribuição.

\begin{itemize}
    \item \emph{Viés}: mede o quanto as previsões do modelo são sistematicamente erradas, independentemente dos dados de treinamento. Um modelo linear aplicado a um problema intrinsecamente quadrático tem alto viés.
    \item \emph{Variância}: mede o quanto as previsões do modelo mudam quando o conjunto de treinamento muda. Um modelo polinomial de grau alto, treinado em amostras diferentes da mesma distribuição, produz previsões muito diferentes (alta variância).
    \item \emph{Ruído irredutível}: a variabilidade inerente dos dados ($y = f(\mathbf{x}) + \varepsilon$), que nenhum modelo pode eliminar.
\end{itemize}

\begin{notabox}{Viés-variância no AR}
O dilema viés-variância já apareceu implicitamente nas notas de AR, na comparação entre Monte Carlo e TD. MC tem baixo viés (usa retornos reais) mas alta variância (o retorno depende de toda a sequência de recompensas ruidosas). TD tem viés introduzido pelo bootstrap (o alvo usa estimativas), mas menor variância (o alvo é mais estável). A escolha entre MC e TD é, portanto, uma escolha no espaço viés-variância.
\end{notabox}

\subsection{Regularização}
\label{subsec:regularizacao}

A regularização é a técnica principal para controlar overfitting sem aumentar o conjunto de dados. A ideia é adicionar um \emph{termo de penalidade} à função de custo que favorece modelos mais simples (com pesos menores):

\[
J_{\text{reg}}(\mathbf{w}) = J(\mathbf{w}) + \lambda \cdot \Omega(\mathbf{w})
\]

onde $\lambda > 0$ é o \emph{hiperparâmetro de regularização} que controla o equilíbrio entre ajuste aos dados e simplicidade do modelo. As duas formas mais comuns são:

\medskip
\noindent\textbf{Regularização $\ell_2$ (Ridge):}
\[
\Omega(\mathbf{w}) = \|\mathbf{w}\|_2^2 = \sum_j w_j^2
\]

Penaliza pesos grandes, forçando-os a permanecer próximos de zero. Todos os pesos são reduzidos, mas raramente chegam a zero exatamente.

\medskip
\noindent\textbf{Regularização $\ell_1$ (Lasso):}
\[
\Omega(\mathbf{w}) = \|\mathbf{w}\|_1 = \sum_j |w_j|
\]

Produz soluções \emph{esparsas}: muitos pesos tornam-se exatamente zero, efetivamente selecionando um subconjunto de características.

% Uma forma intuitiva de visualizar a regularização é imaginar que o otimizador procura a menor perda, mas dentro de uma região permitida para os pesos. A forma dessa região explica por que $\ell_2$ apenas encolhe pesos, enquanto $\ell_1$ tende a zerar alguns deles, como mostra a Figura~\ref{fig:regularizacao}.
Uma forma intuitiva de visualizar a regularização é interpretar a restrição sobre os pesos como uma região permitida no plano $(w_1,w_2)$. Sem regularização, o otimizador buscaria o centro das curvas de nível da perda, isto é, o ponto de menor erro nos dados de treinamento. Com regularização, porém, a solução deve permanecer dentro da região permitida. A solução regularizada ocorre no primeiro ponto em que uma curva de nível da perda toca essa região, como mostra a Figura~\ref{fig:regularizacao}. Na regularização $\ell_2$, a região permitida é circular: sua borda é suave, sem cantos alinhados aos eixos. Por isso, o ponto de contato costuma ter coordenadas pequenas, mas raramente exatamente nulas. Já na regularização $\ell_1$, a região permitida tem forma de losango, com vértices sobre os eixos $w_1=0$ e $w_2=0$. Como esses vértices são pontos salientes da região, é frequente que uma curva de nível toque o losango justamente em um deles. Quando isso acontece, uma das coordenadas do vetor de pesos fica exatamente igual a zero. Essa é a origem geométrica da esparsidade induzida pela regularização $\ell_1$.

\begin{figure}[ht]
\centering
\begin{tikzpicture}[
    lbl/.style={font=\scriptsize, align=center}
]
% L2
\begin{scope}[xshift=0cm]
\draw[->] (-2.5,0) -- (2.5,0) node[right,font=\tiny] {$w_1$};
\draw[->] (0,-2.5) -- (0,2.5) node[above,font=\tiny] {$w_2$};
% L2 ball
\draw[thick, agentcolor] (0,0) circle[radius=1.5];
% Contornos de custo (elipses)
\foreach \r in {0.8, 1.3, 1.9} {
    \draw[gray!40, domain=0:360, samples=60]
        plot ({1.8*\r*cos(\x)-0.9}, {0.7*\r*sin(\x)+0.9});
}
\node[circle, fill=mlred, minimum size=4pt, inner sep=0] at (-0.47, 1.43) {};
\node[lbl, text=mlred] at (-1.3, 2.0) {$\mathbf{w}^*_{\ell_2}$};
\node[lbl, text=agentcolor] at (0,-3.0) {$\ell_2$: região circular};
\node[lbl, text=black!60, font=\tiny] at (0,-3.7) {pesos pequenos, não exatamente 0};
\end{scope}

% L1
\begin{scope}[xshift=6.5cm]
\draw[->] (-2.5,0) -- (2.5,0) node[right,font=\tiny] {$w_1$};
\draw[->] (0,-2.5) -- (0,2.5) node[above,font=\tiny] {$w_2$};
% L1 ball (losango)
\draw[thick, agentcolor] (1.5,0) -- (0,1.5) -- (-1.5,0) -- (0,-1.5) -- cycle;
% Contornos de custo (elipses)
\foreach \r in {0.8, 1.3, 1.9} {
    \draw[gray!40, domain=0:360, samples=60]
        plot ({1.8*\r*cos(\x)-0.9}, {0.7*\r*sin(\x)+0.9});
}
\node[circle, fill=mlred, minimum size=4pt, inner sep=0] at (0, 1.5) {};
\node[lbl, text=mlred] at (0.6, 2.0) {$\mathbf{w}^*_{\ell_1}$};
\node[lbl, text=agentcolor] at (0,-3.0) {$\ell_1$: região em losango};
\node[lbl, text=black!60, font=\tiny] at (0,-3.7) {soluções esparsas ($w_1 = 0$)};
\end{scope}
\end{tikzpicture}
% \caption{Intuição geométrica da regularização $\ell_2$ e $\ell_1$. A solução regularizada é o ponto de menor custo (curvas de nível) que toca a região permitida para os pesos. Na $\ell_2$ (bola circular), a solução se desloca para um ponto arbitrário da borda. Na $\ell_1$ (losango), as soluções tendem a tocar os vértices, onde uma das componentes é zero, daí a esparsidade.}
\caption{Intuição geométrica da regularização $\ell_2$ e $\ell_1$. As curvas cinzas representam níveis de mesma perda, e a região azul representa os pesos permitidos pela regularização. A solução regularizada é o ponto em que a menor curva de nível possível toca essa região. Na $\ell_2$, a borda circular tende a produzir pesos menores, mas não exatamente nulos. Na $\ell_1$, o losango possui vértices sobre os eixos; quando a solução toca um desses vértices, uma coordenada do vetor de pesos torna-se exatamente zero, produzindo uma solução esparsa.}
\label{fig:regularizacao}
\end{figure}

\subsubsection{Exemplo resolvido}

Suponha um modelo linear com dois pesos $w_1 = 3{,}0$ e $w_2 = 0{,}1$, e custo de treinamento $J = 1{,}5$, com $\lambda = 0{,}5$.

\paragraph{Custo regularizado $\ell_2$:}
\[
J_{\ell_2} = 1{,}5 + 0{,}5 \times (3{,}0^2 + 0{,}1^2) = 1{,}5 + 0{,}5 \times 9{,}01 = 1{,}5 + 4{,}505 = 6{,}005
\]

\paragraph{Custo regularizado $\ell_1$:}
\[
J_{\ell_1} = 1{,}5 + 0{,}5 \times (|3{,}0| + |0{,}1|) = 1{,}5 + 0{,}5 \times 3{,}1 = 1{,}5 + 1{,}55 = 3{,}05
\]

\paragraph{Interpretação:} a regularização penaliza $w_1 = 3{,}0$ muito mais do que $w_2 = 0{,}1$. O otimizador é incentivado a reduzir $w_1$, a menos que mantê-lo grande seja fortemente justificado pelos dados. Isso previne que um único peso domine o modelo e cause overfitting.

\subsection{Validação cruzada}
\label{subsec:cross-val}

Nos esquemas anteriores, vimos duas formas comuns de separar dados: \emph{train/test} e \emph{train/val/test}. No segundo caso, o conjunto de validação é usado durante o desenvolvimento do modelo: escolhemos hiperparâmetros, comparamos arquiteturas e decidimos quando parar o treinamento com base no desempenho nesse subconjunto. O conjunto de teste, por sua vez, permanece intocado até a avaliação final.

Essa estratégia funciona bem quando há dados suficientes. Mas, em conjuntos pequenos, reservar uma parte fixa dos dados apenas para validação pode ser caro: o modelo passa a ser treinado com menos exemplos, e a estimativa de desempenho pode depender muito de quais exemplos caíram no subconjunto de validação. A validação cruzada surge exatamente para reduzir esse problema.

A ideia é substituir o conjunto de validação fixo por várias validações sucessivas. Em vez de separar uma única fatia dos dados para validação, o conjunto de treino é dividido em partes, e cada parte será usada como validação uma vez. Assim, todos os exemplos participam do treinamento em algumas iterações e da validação em outra. O resultado é uma estimativa de desempenho menos dependente de uma única divisão dos dados.

Quando o conjunto de dados é pequeno, reservar 15--20\% para validação pode desperdiçar dados preciosos de treinamento. A \emph{validação cruzada} $k$-fold resolve esse problema:
% Quando o conjunto de dados é pequeno, reservar 15--20\% para validação pode desperdiçar dados preciosos de treinamento. A \emph{validação cruzada} $k$-fold resolve esse problema:

\begin{enumerate}
    \item Divida o conjunto de treino em $k$ partes (\textit{folds}) de tamanho igual.
    \item Para $i = 1, \ldots, k$: treine o modelo em todos os folds exceto o $i$-ésimo; avalie no fold $i$.
    \item O erro de validação estimado é a média dos $k$ erros obtidos.
\end{enumerate}

A Figura~\ref{fig:cross-val} ilustra o funcionamento da validação cruzada 5-fold. O conjunto de treino é dividido em cinco partes, e o processo é repetido cinco vezes. Em cada iteração, uma parte diferente é usada como validação, enquanto as quatro restantes são usadas para treinamento. Ao final, obtêm-se cinco estimativas de erro, uma para cada escolha do fold de validação, e o desempenho final é calculado pela média desses erros. A figura também reforça que essa alternância ocorre apenas dentro dos dados disponíveis para treinamento e validação; o conjunto de teste permanece separado e não participa desse processo.

\begin{figure}[ht]
\centering
\begin{tikzpicture}[xscale=0.95]
% Helper sizes
\def\fw{2.1}  % fold width
\def\fh{0.52} % fold height
\def\ys{0.72} % y spacing between rows

% Row 1: val=fold 0
\node[font=\scriptsize, text=black!70] at (-0.5, 0*\ys+\fh/2) {iter 1};
\fill[alertbg]  (0*\fw,0*\ys) rectangle (1*\fw,0*\ys+\fh); \draw[alertframe,thick] (0*\fw,0*\ys) rectangle (1*\fw,0*\ys+\fh);
\fill[statecolor](1*\fw,0*\ys) rectangle (2*\fw,0*\ys+\fh); \draw[agentcolor,thick] (1*\fw,0*\ys) rectangle (2*\fw,0*\ys+\fh);
\fill[statecolor](2*\fw,0*\ys) rectangle (3*\fw,0*\ys+\fh); \draw[agentcolor,thick] (2*\fw,0*\ys) rectangle (3*\fw,0*\ys+\fh);
\fill[statecolor](3*\fw,0*\ys) rectangle (4*\fw,0*\ys+\fh); \draw[agentcolor,thick] (3*\fw,0*\ys) rectangle (4*\fw,0*\ys+\fh);
\fill[statecolor](4*\fw,0*\ys) rectangle (5*\fw,0*\ys+\fh); \draw[agentcolor,thick] (4*\fw,0*\ys) rectangle (5*\fw,0*\ys+\fh);
\node[font=\scriptsize\bfseries, text=alertframe] at (0.5*\fw,0*\ys+\fh/2) {val};
\foreach \c in {1,2,3,4}{\node[font=\tiny, text=agentcolor] at (\c*\fw+0.5*\fw,0*\ys+\fh/2) {treino};}

% Row 2: val=fold 1
\node[font=\scriptsize, text=black!70] at (-0.5, 1*\ys+\fh/2) {iter 2};
\fill[statecolor](0*\fw,1*\ys) rectangle (1*\fw,1*\ys+\fh); \draw[agentcolor,thick] (0*\fw,1*\ys) rectangle (1*\fw,1*\ys+\fh);
\fill[alertbg]   (1*\fw,1*\ys) rectangle (2*\fw,1*\ys+\fh); \draw[alertframe,thick] (1*\fw,1*\ys) rectangle (2*\fw,1*\ys+\fh);
\fill[statecolor](2*\fw,1*\ys) rectangle (3*\fw,1*\ys+\fh); \draw[agentcolor,thick] (2*\fw,1*\ys) rectangle (3*\fw,1*\ys+\fh);
\fill[statecolor](3*\fw,1*\ys) rectangle (4*\fw,1*\ys+\fh); \draw[agentcolor,thick] (3*\fw,1*\ys) rectangle (4*\fw,1*\ys+\fh);
\fill[statecolor](4*\fw,1*\ys) rectangle (5*\fw,1*\ys+\fh); \draw[agentcolor,thick] (4*\fw,1*\ys) rectangle (5*\fw,1*\ys+\fh);
\node[font=\tiny, text=agentcolor] at (0.5*\fw,1*\ys+\fh/2) {treino};
\node[font=\scriptsize\bfseries, text=alertframe] at (1.5*\fw,1*\ys+\fh/2) {val};
\foreach \c in {2,3,4}{\node[font=\tiny, text=agentcolor] at (\c*\fw+0.5*\fw,1*\ys+\fh/2) {treino};}

% Row 3: val=fold 2
\node[font=\scriptsize, text=black!70] at (-0.5, 2*\ys+\fh/2) {iter 3};
\fill[statecolor](0*\fw,2*\ys) rectangle (1*\fw,2*\ys+\fh); \draw[agentcolor,thick] (0*\fw,2*\ys) rectangle (1*\fw,2*\ys+\fh);
\fill[statecolor](1*\fw,2*\ys) rectangle (2*\fw,2*\ys+\fh); \draw[agentcolor,thick] (1*\fw,2*\ys) rectangle (2*\fw,2*\ys+\fh);
\fill[alertbg]   (2*\fw,2*\ys) rectangle (3*\fw,2*\ys+\fh); \draw[alertframe,thick] (2*\fw,2*\ys) rectangle (3*\fw,2*\ys+\fh);
\fill[statecolor](3*\fw,2*\ys) rectangle (4*\fw,2*\ys+\fh); \draw[agentcolor,thick] (3*\fw,2*\ys) rectangle (4*\fw,2*\ys+\fh);
\fill[statecolor](4*\fw,2*\ys) rectangle (5*\fw,2*\ys+\fh); \draw[agentcolor,thick] (4*\fw,2*\ys) rectangle (5*\fw,2*\ys+\fh);
\foreach \c in {0,1}{\node[font=\tiny, text=agentcolor] at (\c*\fw+0.5*\fw,2*\ys+\fh/2) {treino};}
\node[font=\scriptsize\bfseries, text=alertframe] at (2.5*\fw,2*\ys+\fh/2) {val};
\foreach \c in {3,4}{\node[font=\tiny, text=agentcolor] at (\c*\fw+0.5*\fw,2*\ys+\fh/2) {treino};}

% Row 4: val=fold 3
\node[font=\scriptsize, text=black!70] at (-0.5, 3*\ys+\fh/2) {iter 4};
\fill[statecolor](0*\fw,3*\ys) rectangle (1*\fw,3*\ys+\fh); \draw[agentcolor,thick] (0*\fw,3*\ys) rectangle (1*\fw,3*\ys+\fh);
\fill[statecolor](1*\fw,3*\ys) rectangle (2*\fw,3*\ys+\fh); \draw[agentcolor,thick] (1*\fw,3*\ys) rectangle (2*\fw,3*\ys+\fh);
\fill[statecolor](2*\fw,3*\ys) rectangle (3*\fw,3*\ys+\fh); \draw[agentcolor,thick] (2*\fw,3*\ys) rectangle (3*\fw,3*\ys+\fh);
\fill[alertbg]   (3*\fw,3*\ys) rectangle (4*\fw,3*\ys+\fh); \draw[alertframe,thick] (3*\fw,3*\ys) rectangle (4*\fw,3*\ys+\fh);
\fill[statecolor](4*\fw,3*\ys) rectangle (5*\fw,3*\ys+\fh); \draw[agentcolor,thick] (4*\fw,3*\ys) rectangle (5*\fw,3*\ys+\fh);
\foreach \c in {0,1,2}{\node[font=\tiny, text=agentcolor] at (\c*\fw+0.5*\fw,3*\ys+\fh/2) {treino};}
\node[font=\scriptsize\bfseries, text=alertframe] at (3.5*\fw,3*\ys+\fh/2) {val};
\node[font=\tiny, text=agentcolor] at (4.5*\fw,3*\ys+\fh/2) {treino};

% Row 5: val=fold 4
\node[font=\scriptsize, text=black!70] at (-0.5, 4*\ys+\fh/2) {iter 5};
\fill[statecolor](0*\fw,4*\ys) rectangle (1*\fw,4*\ys+\fh); \draw[agentcolor,thick] (0*\fw,4*\ys) rectangle (1*\fw,4*\ys+\fh);
\fill[statecolor](1*\fw,4*\ys) rectangle (2*\fw,4*\ys+\fh); \draw[agentcolor,thick] (1*\fw,4*\ys) rectangle (2*\fw,4*\ys+\fh);
\fill[statecolor](2*\fw,4*\ys) rectangle (3*\fw,4*\ys+\fh); \draw[agentcolor,thick] (2*\fw,4*\ys) rectangle (3*\fw,4*\ys+\fh);
\fill[statecolor](3*\fw,4*\ys) rectangle (4*\fw,4*\ys+\fh); \draw[agentcolor,thick] (3*\fw,4*\ys) rectangle (4*\fw,4*\ys+\fh);
\fill[alertbg]   (4*\fw,4*\ys) rectangle (5*\fw,4*\ys+\fh); \draw[alertframe,thick] (4*\fw,4*\ys) rectangle (5*\fw,4*\ys+\fh);
\foreach \c in {0,1,2,3}{\node[font=\tiny, text=agentcolor] at (\c*\fw+0.5*\fw,4*\ys+\fh/2) {treino};}
\node[font=\scriptsize\bfseries, text=alertframe] at (4.5*\fw,4*\ys+\fh/2) {val};

\node[font=\scriptsize, text=black!60, align=center] at (2.5*\fw, -0.55)
{5-fold: cada iteração usa um fold diferente como validação. Erro final = média dos 5 erros.};
\end{tikzpicture}
\caption{Validação cruzada 5-fold. A cada iteração, um fold diferente é usado para validação e os restantes para treinamento. 
O erro estimado é a média dos 5 erros de validação. Ao longo das 5 iterações, cada exemplo é usado para validação uma vez e para treinamento nas demais iterações.
O conjunto de teste permanece completamente separado e não participa da validação cruzada.}
\label{fig:cross-val}
\end{figure}

\begin{alertabox}{Validação cruzada não substitui o conjunto de teste}
A validação cruzada substitui a separação treino/validação durante a escolha de modelos e hiperparâmetros, permitindo usar os dados de forma mais eficiente. No entanto, ela não substitui o conjunto de teste. O teste continua reservado para a avaliação final, depois que todas as escolhas de modelo e hiperparâmetros foram feitas. Aplicar validação cruzada ao conjunto de teste é um erro equivalente ao vazamento de dados.
\end{alertabox}

\subsection{Métricas de avaliação}
\label{subsec:metricas}

A escolha da métrica de avaliação depende da tarefa e dos requisitos do domínio.

\paragraph{Regressão.}
\begin{itemize}[nosep]
    \item \emph{MSE} (Erro Quadrático Médio): $\frac{1}{n}\sum (y_i - \hat{y}_i)^2$. Penaliza erros grandes.
    \item \emph{RMSE}: $\sqrt{\text{MSE}}$. Na mesma unidade da variável-alvo.
    \item \emph{MAE} (Erro Absoluto Médio): $\frac{1}{n}\sum |y_i - \hat{y}_i|$. Mais robusto a outliers.
    \item \emph{$R^2$} (coeficiente de determinação): fração da variância explicada pelo modelo. $R^2 = 1$ indica ajuste perfeito; $R^2 = 0$, não melhor que prever a média.
\end{itemize}

\paragraph{Classificação.}

A métrica mais comum é a \emph{acurácia}: fração de exemplos classificados corretamente. Mas ela pode ser enganosa em dados desbalanceados. Se 95\% dos exemplos são da classe $0$, um modelo que sempre prevê $0$ tem acurácia 95\% mas é inútil.

Para analisar melhor o desempenho, usamos a \emph{matriz de confusão} (Tabela~\ref{tab:matriz-confusao}) e métricas derivadas. A distinção entre falso positivo e falso negativo é essencial porque os dois tipos de erro podem ter custos muito diferentes dependendo do domínio.

\begin{table}[ht]
\centering
\caption{Matriz de confusão para classificação binária.}
\label{tab:matriz-confusao}
\renewcommand{\arraystretch}{1.3}
\begin{tabular}{llcc}
\toprule
 & & \multicolumn{2}{c}{\textbf{Previsto}} \\
 & & Positivo & Negativo \\
\midrule
\multirow{2}{*}{\textbf{Real}} & Positivo & VP (Verdadeiro Positivo) & FN (Falso Negativo) \\
                                & Negativo & FP (Falso Positivo) & VN (Verdadeiro Negativo) \\
\bottomrule
\end{tabular}
\end{table}

\begin{itemize}[nosep]
    \item \textbf{Precisão}: $\frac{\text{VP}}{\text{VP}+\text{FP}}$ — dos que foram classificados como positivos, quantos realmente são?
    \item \textbf{Revocação} (\textit{recall}): $\frac{\text{VP}}{\text{VP}+\text{FN}}$ — dos que são positivos, quantos o modelo encontrou?
    \item \textbf{F1-score}: média harmônica de precisão e revocação: $\frac{2 \cdot \text{Precisão} \cdot \text{Revocação}}{\text{Precisão}+\text{Revocação}}$.
\end{itemize}

\begin{notabox}{Precisão \textit{versus} revocação: o dilema do médico}
Em diagnóstico médico, um falso negativo (não detectar uma doença grave) é geralmente muito pior do que um falso positivo (suspeitar de uma doença que não existe). Nesse contexto, maximizar a revocação é prioritário, mesmo que isso reduza a precisão. Já em filtragem de spam, um falso positivo (deletar um e-mail legítimo) pode ser pior do que um falso negativo (deixar passar um spam). A escolha da métrica deve refletir as consequências reais dos erros no domínio.
\end{notabox}

\begin{checkpointbox}{Micro-checkpoint}
Um modelo de AM para detecção de fraude em transações bancárias tem 99{,}8\% de acurácia, mas a equipe de segurança está insatisfeita. Explique por que a acurácia pode ser uma métrica enganosa neste contexto. Qual combinação de métricas seria mais informativa?
\end{checkpointbox}

\begin{perguntabox}
Um engenheiro de ML treina três versões de um modelo de regressão usando características polinomiais de graus 1, 5 e 15. O modelo de grau 1 obtém MSE de treino $= 120$ e MSE de validação $= 125$. O de grau 5 obtém MSE de treino $= 45$ e MSE de validação $= 52$. O de grau 15 obtém MSE de treino $= 2$ e MSE de validação $= 380$. \textbf{(a)} Qual modelo apresenta underfitting? Qual apresenta overfitting? \textbf{(b)} Qual você selecionaria para implantação e por quê? \textbf{(c)} Se você quisesse melhorar o modelo de grau 15 sem reduzir o grau, que estratégia usaria?
\end{perguntabox}

%===========================================
\section{Árvores de decisão}
\label{subsec:arvores}

Modelos lineares representam fronteiras de decisão como hiperplanos. 
As \emph{árvores de decisão} adotam outra estratégia: em vez de calcular uma única pontuação a partir de todas as variáveis, elas conduzem o exemplo $\mathbf{x}$ por uma sequência de perguntas simples. Cada pergunta testa uma característica $x_j$ do exemplo. Cada resposta leva a uma nova pergunta ou a uma decisão final.

%Modelos lineares representam fronteiras de decisão como hiperplanos. As \emph{árvores de decisão} adotam uma abordagem radicalmente diferente: particionam o espaço de entradas em regiões hierárquicas por meio de uma sequência de testes simples (geralmente do tipo $x_j \leq \tau$), e fazem uma previsão constante em cada região folha.

Antes de examinar uma árvore formalmente, considere um problema simples: classificar um e-mail como \emph{spam} ou \emph{legítimo}. Um classificador linear tentaria combinar várias características em uma única pontuação. Uma árvore de decisão segue outra lógica: ela faz uma pergunta por vez. Por exemplo, suponha que usamos três características de um e-mail:

\begin{itemize}
    \item quantas vezes aparece a palavra ``grátis'';
    \item quantos links externos o e-mail contém;
    \item se o assunto está escrito em letras maiúsculas.
\end{itemize}

Uma possível regra de classificação seria a seguinte. Primeiro, verificamos se a palavra ``grátis'' aparece pelo menos três vezes. Se sim, olhamos a quantidade de links externos. Se houver muitos links, classificamos como spam; caso contrário, como legítimo. Se a palavra ``grátis'' não aparece com tanta frequência, verificamos se o assunto está em maiúsculas. Dependendo dessa resposta, chegamos novamente a uma das duas previsões possíveis.

Essa sequência de perguntas pode ser representada como uma árvore. Cada nó interno corresponde a um teste sobre uma característica. Cada aresta corresponde a uma resposta possível ao teste. Cada folha contém a decisão final.

A Figura~\ref{fig:arvore} apresenta uma árvore de decisão para classificar e-mails como spam ou legítimos. A classificação é feita por uma sequência de perguntas simples sobre atributos do e-mail. Primeiro, verifica-se se a palavra ``grátis'' aparece pelo menos três vezes. Dependendo da resposta, o exemplo segue por um dos ramos da árvore e passa por novos testes, como a quantidade de links externos ou o uso de assunto em letras maiúsculas. O caminho termina em uma folha, que contém a previsão final. Essa estrutura torna o modelo interpretável: para justificar uma decisão, basta seguir o caminho percorrido da raiz até a folha correspondente. No exemplo, um e-mail pode ser marcado como spam não por uma pontuação abstrata, mas porque satisfez uma combinação explícita de condições.

\begin{figure}[ht]
\centering
\begin{tikzpicture}[
    inode/.style={draw=agentcolor, fill=statecolor, rounded corners=3pt, thick,
                  minimum width=3.2cm, minimum height=0.8cm,
                  align=center, font=\small},
    lnode/.style={draw=envcolor, fill=workedbg, rounded corners=3pt, thick,
                  minimum width=2.4cm, minimum height=0.8cm,
                  align=center, font=\small},
    lnodebad/.style={draw=mlred, fill=termneg, rounded corners=3pt, thick,
                  minimum width=2.4cm, minimum height=0.8cm,
                  align=center, font=\small},
    arr/.style={-{Stealth[length=5pt]}, thick, black!60},
    lbl/.style={font=\scriptsize, fill=white, inner sep=1pt}
]

\node[inode] (r)  at (5,  0)   {``grátis'' $\geq 3$?};
\node[inode] (l1) at (2, -1.8) {links externos $> 2$?};
\node[inode] (r1) at (8, -1.8) {assunto maiúsculo?};

\node[lnodebad] (ll1) at (0.5, -3.6) {\faEnvelope~spam};
\node[lnode]    (lr1) at (3.5, -3.6) {\faCheck~legítimo};
\node[lnodebad] (rl1) at (6.5, -3.6) {\faEnvelope~spam};
\node[lnode]    (rr1) at (9.5, -3.6) {\faCheck~legítimo};

\draw[arr] (r) -- (l1) node[lbl, midway, above left] {sim};
\draw[arr] (r) -- (r1) node[lbl, midway, above right] {não};
\draw[arr] (l1) -- (ll1) node[lbl, midway, left]  {sim};
\draw[arr] (l1) -- (lr1) node[lbl, midway, right] {não};
\draw[arr] (r1) -- (rl1) node[lbl, midway, left]  {sim};
\draw[arr] (r1) -- (rr1) node[lbl, midway, right] {não};

\end{tikzpicture}
\caption{Exemplo de árvore de decisão para classificação de spam. Cada nó interno testa uma característica; cada folha produz uma previsão. A interpretabilidade é uma vantagem central: é fácil explicar por que um e-mail foi classificado como spam.}
\label{fig:arvore}
\end{figure}

De forma mais geral, uma árvore de decisão divide o espaço de entradas em regiões sucessivamente menores. Cada nó interno aplica um teste, geralmente da forma $x_j \leq \tau$, em que $x_j$ representa a $j$-ésima característica do exemplo e $\tau$ é um limiar escolhido pelo algoritmo. Cada resposta ao teste define um ramo da árvore. Ao final do caminho, chega-se a uma folha, isto é, um nó terminal que contém a previsão do modelo. Como todos os exemplos que chegam à mesma folha recebem a mesma saída, dizemos que a árvore faz uma previsão constante em cada região associada a uma folha.

A principal vantagem das árvores de decisão é a \emph{interpretabilidade}: cada decisão pode ser rastreada por um caminho explícito na árvore. A principal desvantagem é que, sem controle de profundidade, árvores profundas tendem a sofrer overfitting severo.

\subsection{Construção de árvores de decisão para classificação}
\label{subsec:entropia}

Nesta seção, consideramos árvores de decisão para \emph{classificação}. Nesse cenário, cada exemplo pertence a uma classe, como \emph{spam} ou \emph{legítimo}. Ao construir a árvore, queremos escolher testes que separem os exemplos em subconjuntos cada vez mais homogêneos em relação à classe.

\subsubsection{Entropia e ganho de informação}

O problema central na construção de uma árvore é: \emph{qual atributo usar em cada nó de divisão?} Em um nó, temos um conjunto de exemplos de treinamento. Se esses exemplos misturam várias classes, o nó é impuro. Se quase todos pertencem à mesma classe, o nó é mais puro. O critério de \emph{ganho de informação} mede quanto uma divisão reduz essa impureza.

A construção da árvore é um procedimento \emph{guloso} (\textit{greedy}): em cada nó, escolhe-se a melhor divisão local segundo o critério de ganho de informação. Isso não garante a árvore globalmente ótima, mas produz um algoritmo eficiente e interpretável.

Para um conjunto de exemplos $S$, seja $p_k$ a proporção de exemplos de $S$ pertencentes à classe $k$. A entropia de $S$ é definida por:

\[
H(S) = -\sum_k p_k \log_2 p_k
\]

A entropia mede a impureza (ou desordem) do conjunto. Ela é máxima quando as classes estão igualmente distribuídas ($H = \log_2 K$ para $K$ classes) e mínima ($H = 0$) quando todos os exemplos pertencem à mesma classe.

A Figura~\ref{fig:entropia-binaria} ilustra o comportamento da entropia no caso binário ($k=2$). Se denotamos por $p$ a proporção de exemplos de uma das classes em um nó, então a outra classe tem proporção $1-p$. Observa-se que a entropia é igual a zero quando $p=0$ ou $p=1$, isto é, quando todos os exemplos pertencem à mesma classe e o nó é puro. À medida que as classes se tornam mais equilibradas, a entropia aumenta, atingindo seu valor máximo em $p=0.5$. Isso confirma a intuição de que a entropia mede impureza: ela é baixa em nós homogêneos e alta em nós com forte mistura de classes.

\begin{figure}[ht]
\centering
\begin{tikzpicture}
\begin{axis}[
    width=10cm,
    height=6.2cm,
    xmin=0, xmax=1,
    ymin=0, ymax=1.1,
    domain=0.001:0.999,
    samples=200,
    axis lines=left,
    xlabel={$p$},
    ylabel={$H(p)$},
    xlabel style={font=\small},
    ylabel style={font=\small},
    xtick={0,0.25,0.5,0.75,1},
    ytick={0,0.5,1},
    grid=both,
    grid style={dashed, gray!30},
    enlargelimits=false,
    clip=false
]
\addplot[thick, agentcolor]
    {-x*ln(x)/ln(2) - (1-x)*ln(1-x)/ln(2)};

\addplot[dashed, black!60] coordinates {(0.5,0) (0.5,1)};
\addplot[dashed, black!60] coordinates {(0,1) (0.5,1)};

\addplot[only marks, mark=*, mark size=2pt] coordinates {(0,0) (0.5,1) (1,0)};

\node[font=\scriptsize, anchor=south west] at (axis cs:0.52,1.0) {máximo em $p=0.5$};
\node[font=\scriptsize, anchor=north west] at (axis cs:0.02,0.03) {nó puro};
\node[font=\scriptsize, anchor=north east] at (axis cs:0.98,0.03) {nó puro};
\end{axis}
\end{tikzpicture}
\caption{Entropia no caso binário, em que uma classe tem proporção $p$ e a outra tem proporção $1-p$. A entropia é mínima quando o nó é puro ($p=0$ ou $p=1$) e máxima quando as duas classes estão igualmente representadas ($p=0.5$).}
\label{fig:entropia-binaria}
\end{figure}

A entropia, sozinha, apenas descreve o grau de mistura das classes em um nó. Para construir a árvore, precisamos dar um passo a mais: avaliar se uma possível pergunta reduz essa mistura. Uma boa pergunta é aquela que pega um conjunto inicialmente impuro e o separa em subconjuntos mais puros.

Por exemplo, no problema de spam, a pergunta ``a palavra grátis aparece pelo menos três vezes?'' divide os e-mails em dois subconjuntos: os que respondem ``sim'' e os que respondem ``não''. Se, depois dessa divisão, cada subconjunto ficar mais homogêneo em termos de classe, então essa pergunta foi útil. O \emph{ganho de informação} mede exatamente essa utilidade: quanto a entropia diminui depois da divisão.

O \emph{ganho de informação} de um atributo $A$ com valores possíveis $\{v_1, v_2, \ldots\}$ é:

\[
\text{IG}(S, A) = H(S) - \sum_v \frac{|S_v|}{|S|} H(S_v)
\]

onde $S_v$ é o subconjunto de $S$ em que $A = v$. O atributo escolhido para divisão é aquele que maximiza $\text{IG}(S, A)$.

\subsubsection{Testes candidatos e construção recursiva da árvore}

Até aqui, falamos em escolher um atributo para dividir um nó. Essa descrição é
útil, mas ainda incompleta. Na prática, o algoritmo não escolhe apenas um
atributo isolado; ele escolhe um \emph{teste}. Um teste define uma pergunta que
separa os exemplos de um nó em subconjuntos.

Quando o atributo é categórico, a divisão costuma ser direta. Por exemplo, o
atributo ``assunto em letras maiúsculas'' já sugere uma pergunta binária:
``o assunto está em letras maiúsculas?''. As respostas possíveis são ``sim'' e
``não''.

Quando o atributo é numérico, o algoritmo também precisa escolher um limiar. Se
$x_j$ representa a $j$-ésima característica de um exemplo, um teste típico tem a
forma

\[
x_j \leq \tau,
\]

em que $\tau$ é o limiar usado para separar os exemplos. No problema de spam, se
$x_j$ representa o número de vezes que a palavra ``grátis'' aparece no e-mail,
um teste possível seria ``grátis'' $\geq 3$. Outros testes também poderiam ser
avaliados, como ``grátis'' $\geq 1$, ``grátis'' $\geq 2$ ou ``grátis'' $\geq 4$.

O limiar não é escolhido arbitrariamente. Em uma implementação real, o algoritmo
testa vários limiares candidatos para cada atributo numérico, calcula o ganho de
informação produzido por cada teste e escolhe aquele que gera a maior redução de
impureza no nó atual. Assim, em atributos numéricos, a escolha feita pelo
algoritmo envolve o par

\[
(\text{atributo}, \text{limiar}).
\]

Depois que o melhor teste é escolhido, o nó é dividido em filhos. Cada filho recebe apenas os exemplos que satisfazem a resposta correspondente ao teste. O mesmo procedimento é então repetido em cada filho: calcula-se a impureza do novo conjunto, avaliam-se novos testes candidatos e escolhe-se a melhor divisão local.

Essa repetição caracteriza a construção \emph{recursiva} da árvore. A raiz é dividida primeiro. Em seguida, cada nó filho pode ser dividido novamente. O processo continua até que algum critério de parada seja satisfeito. Há vários critérios que podem ser usados. Por exemplo: o nó já contém exemplos de uma única classe, não há ganho de informação suficiente, a profundidade máxima foi atingida ou há poucos exemplos no nó.

Portanto, uma árvore de decisão é construída por uma sequência de decisões locais. Em cada nó, o algoritmo pergunta: ``qual teste torna os filhos mais puros?''. A resposta define a divisão daquele nó. A aplicação recursiva dessa ideia produz a árvore completa.

\subsubsection{Exemplo resolvido}

Considere um conjunto de 10 e-mails: 6 spam ($+$) e 4 legítimos ($-$). 
Suponha que o algoritmo esteja comparando dois testes candidatos para dividir o nó raiz:

\begin{itemize}[nosep]
    \item Teste A: ``grátis'' $\geq 3$, que divide os exemplos em 
    $S_{\text{sim}} = \{5{+}, 1{-}\}$ e 
    $S_{\text{não}} = \{1{+}, 3{-}\}$.
    
    \item Teste B: links $> 2$, que divide os exemplos em 
    $S_{\text{sim}} = \{3{+}, 3{-}\}$ e 
    $S_{\text{não}} = \{3{+}, 1{-}\}$.
\end{itemize}

\paragraph{Passo 1: calcular a entropia do conjunto raiz.}
\[
H(S) = -\frac{6}{10}\log_2\frac{6}{10} - \frac{4}{10}\log_2\frac{4}{10}
= -0{,}6 \times (-0{,}737) - 0{,}4 \times (-1{,}322)
= 0{,}442 + 0{,}529 = 0{,}971
\]

\paragraph{Passo 2: calcular IG do Atributo A.}
\[
H(S_{\text{sim}}) = -\frac{5}{6}\log_2\frac{5}{6} - \frac{1}{6}\log_2\frac{1}{6} \approx 0{,}650
\]
\[
H(S_{\text{não}}) = -\frac{1}{4}\log_2\frac{1}{4} - \frac{3}{4}\log_2\frac{3}{4} \approx 0{,}811
\]
\[
\text{IG}(S, A) = 0{,}971 - \left(\frac{6}{10} \times 0{,}650 + \frac{4}{10} \times 0{,}811\right)
= 0{,}971 - (0{,}390 + 0{,}324) = 0{,}971 - 0{,}714 = 0{,}257
\]

\paragraph{Passo 3: calcular IG do Atributo B.}
\[
H(S_{\text{sim}}) = -\frac{3}{6}\log_2\frac{3}{6} - \frac{3}{6}\log_2\frac{3}{6} = 1{,}0
\]
\[
H(S_{\text{não}}) = -\frac{3}{4}\log_2\frac{3}{4} - \frac{1}{4}\log_2\frac{1}{4} \approx 0{,}811
\]
\[
\text{IG}(S, B) = 0{,}971 - \left(\frac{6}{10} \times 1{,}0 + \frac{4}{10} \times 0{,}811\right)
= 0{,}971 - (0{,}600 + 0{,}324) = 0{,}971 - 0{,}924 = 0{,}047
\]

\paragraph{Decisão:} o Atributo A ($\text{IG} = 0{,}257$) é muito melhor que o B ($\text{IG} = 0{,}047$). O algoritmo escolhe A para dividir o nó raiz. Intuitivamente, A cria subconjuntos muito mais puros (mais homogêneos em termos de classe).

Neste exemplo, estamos comparando testes completos, isto é, pares formados por um atributo e um limiar. Em uma árvore real, outros limiares para os mesmos atributos também poderiam ser avaliados.

\subsection{Controle de complexidade e poda}
\label{subsec:poda}

O procedimento de construção de uma árvore é recursivo: em cada nó, o algoritmo procura o teste que produz a maior redução de impureza e divide os exemplos em subconjuntos. Se esse processo continuar sem restrições, a árvore pode crescer até representar detalhes muito específicos do conjunto de treinamento.

No limite, uma árvore pode criar folhas quase puras contendo poucos exemplos, ou até mesmo um único exemplo. Isso parece bom do ponto de vista do erro de treinamento, pois a árvore passa a memorizar muitos casos particulares. Porém, essa memorização geralmente não se transfere para novos dados. A árvore se torna muito profunda, muito sensível a pequenas variações nos dados e tende a sofrer
\emph{overfitting}.

Por isso, árvores de decisão precisam de mecanismos de controle de complexidade. Esses mecanismos cumprem papel semelhante ao da regularização em outros modelos: eles impedem que o modelo se ajuste demais ao conjunto de treinamento. Em árvores, esse controle é feito principalmente por critérios de parada e por poda.

Há duas estratégias principais.

\begin{itemize}
    \item \emph{Pré-poda} (\textit{pre-pruning}): interrompe o crescimento da árvore durante a construção. Um nó deixa de ser dividido quando a divisão não parece suficientemente útil ou quando há poucos dados para justificar uma nova ramificação. Exemplos de critérios são: ganho de informação menor que um limiar, profundidade máxima atingida ou número de exemplos no nó
    menor que $m_{\min}$.

    \item \emph{Pós-poda} (\textit{post-pruning}): primeiro constrói uma árvore grande, possivelmente ajustada demais ao treino, e depois remove ramos que não melhoram o desempenho em dados de validação. A ideia é perguntar: ``se substituirmos este ramo inteiro por uma folha, o desempenho fora do treino piora?''. Se não piorar, o ramo pode ser podado.
\end{itemize}

A diferença entre as duas estratégias está no momento da decisão. Na pré-poda, o algoritmo evita criar certos ramos. Na pós-poda, ele cria os ramos primeiro e decide depois quais devem ser removidos. A pré-poda é mais simples e eficiente, mas pode parar cedo demais. A pós-poda costuma ser mais robusta, pois avalia a utilidade dos ramos depois que a árvore completa revelou sua estrutura.

\begin{alertabox}{A instabilidade das árvores de decisão}
Árvores de decisão são modelos de alta variância. Pequenas mudanças no conjunto de treinamento podem alterar os testes escolhidos perto da raiz, e essas mudanças se propagam para toda a estrutura da árvore. Assim, remover ou adicionar poucos exemplos pode produzir árvores bastante diferentes.

Essa instabilidade é uma desvantagem de uma árvore isolada, mas também explica por que árvores são boas candidatas para métodos ensemble. Se uma única árvore é instável, podemos treinar muitas árvores diferentes e combinar suas previsões. A combinação tende a suavizar os erros individuais e produzir um modelo final mais estável e preciso.
\end{alertabox}

\section{Métodos ensemble}
\label{subsec:ensemble}

Um \emph{método ensemble} combina as previsões de múltiplos modelos (chamados de \emph{modelos base} ou \emph{aprendizes fracos}) para produzir uma previsão final mais robusta. A intuição é que modelos diferentes erram em lugares diferentes, e sua combinação cancela esses erros.

\subsection{Bagging: modelos independentes em paralelo}

O \emph{bagging} (\textit{bootstrap aggregating}) reduz a variância treinando $T$ modelos independentes em subconjuntos diferentes dos dados (amostrados com reposição usando o método \textit{bootstrap}) e combinando suas previsões por votação (classificação) ou média (regressão).

\paragraph{Random Forests.}

Random Forests~\citep{breiman2001random} são a realização mais bem-sucedida do bagging com árvores de decisão. A ideia adicional é introduzir aleatoriedade também na seleção de atributos: em cada nó, em vez de considerar todos os $d$ atributos, considera-se apenas um subconjunto aleatório de $\sqrt{d}$ (para classificação) ou $d/3$ (para regressão) atributos candidatos.

Esse mecanismo de dupla aleatorização (amostras e atributos) produz árvores mais diversas e menos correlacionadas entre si, o que resulta em cancelamento mais eficaz dos erros individuais.

A Figura~\ref{fig:random-forest} mostra a ideia central por trás de \textit{bagging} e Random Forests. A partir do conjunto original de dados $\mathcal{D}$, são geradas várias amostras \textit{bootstrap}, isto é, subconjuntos obtidos por sorteio com reposição. Cada amostra é usada para treinar uma árvore diferente, produzindo modelos que veem versões ligeiramente distintas dos dados. Como árvores de decisão tendem a ter alta variância, pequenas mudanças no conjunto de treino podem gerar árvores bastante diferentes. O ensemble aproveita essa instabilidade: em vez de confiar em uma única árvore, combina as previsões de várias delas por votação, em problemas de classificação, ou por média, em problemas de regressão. A Random Forest acrescenta ainda uma segunda fonte de diversidade, selecionando aleatoriamente um subconjunto de atributos em cada nó da árvore. Essa combinação reduz a variância do modelo final e torna a previsão mais robusta.

\begin{figure}[htb]
\centering
\begin{tikzpicture}[
    box/.style={draw, rounded corners=3pt, thick, minimum width=2.2cm,
                minimum height=0.8cm, align=center, font=\scriptsize},
    arr/.style={-{Stealth[length=5pt]}, thick},
    lbl/.style={font=\scriptsize, fill=white, inner sep=1pt}
]
\node[box, fill=statecolor, draw=agentcolor] (D) at (0,0) {Dados $\mathcal{D}$};

\foreach \i/\x in {1/-4, 2/0, 3/4} {
    \node[box, fill=alertbg, draw=alertframe] (B\i) at (\x, -1.5)
        {amostra\\bootstrap $\i$};
    \node[box, fill=workedbg, draw=envcolor] (T\i) at (\x, -3.0)
        {árvore $\i$};
    \node[box, fill=checkpointbg, draw=checkpointframe] (P\i) at (\x, -4.5)
        {previsão $\i$};
    \draw[arr] (D) -- (B\i);
    \draw[arr] (B\i) -- (T\i);
    \draw[arr] (T\i) -- (P\i);
}

\node[box, fill=boxbg, draw=boxblue, minimum width=3.5cm] (final) at (0, -6.0)
    {Votação / Média\\(previsão final)};
\draw[arr] (P1) -- (final);
\draw[arr] (P2) -- (final);
\draw[arr] (P3) -- (final);

\node[font=\scriptsize, text=black!60, align=center] at (0,-8.0)
{Random Forest: adiciona seleção aleatória de atributos em cada nó};
\end{tikzpicture}
\caption{Arquitetura do bagging / Random Forest. Múltiplas árvores são treinadas em amostras bootstrap independentes; suas previsões são combinadas por votação ou média, reduzindo a variância em relação a uma única árvore.}
\label{fig:random-forest}
\end{figure}

\subsection{Boosting: modelos sequenciais que corrigem erros}

O \emph{boosting} adota uma estratégia diferente: em vez de treinar modelos independentes, treina-os \emph{sequencialmente}, onde cada modelo foca nos exemplos que os anteriores erraram. O algoritmo mais clássico é o \emph{AdaBoost}: exemplos classificados incorretamente recebem peso maior na rodada seguinte, forçando o próximo modelo a se concentrar neles.

O \emph{Gradient Boosting} (e sua implementação eficiente, \emph{XGBoost}~\citep{chen2016xgboost}) generaliza essa ideia: cada novo modelo é treinado para corrigir os \emph{resíduos} (erros) do ensemble acumulado até então. O resultado é uma sequência de modelos simples que, somados, produzem um modelo muito poderoso.

A Figura~\ref{fig:boosting} destaca a diferença central entre boosting e bagging. No bagging, os modelos base são treinados em paralelo e depois combinados. No boosting, os modelos são treinados em sequência. O primeiro modelo produz uma aproximação inicial. Em seguida, o algoritmo identifica os exemplos mal previstos, ou os resíduos deixados pelo ensemble atual, e treina um novo modelo para reduzir esses erros. A previsão final é obtida pela combinação acumulada dos modelos. Por isso, boosting costuma ser associado à redução de viés: ele constrói gradualmente um modelo mais expressivo a partir de vários modelos simples.

\begin{figure}[htb]
\centering
\begin{tikzpicture}[
    box/.style={draw, rounded corners=3pt, thick, minimum width=2.4cm,
                minimum height=0.8cm, align=center, font=\scriptsize},
    arr/.style={-{Stealth[length=5pt]}, thick},
    lbl/.style={font=\scriptsize, fill=white, inner sep=1pt}
]
\node[box, fill=statecolor, draw=agentcolor] (D) at (0,0) {Dados $\mathcal{D}$};

\node[box, fill=workedbg, draw=envcolor] (M1) at (-4,-2) {modelo 1};
\node[box, fill=workedbg, draw=envcolor] (M2) at (0,-2) {modelo 2};
\node[box, fill=workedbg, draw=envcolor] (M3) at (4,-2) {modelo 3};

\node[box, fill=alertbg, draw=alertframe] (E1) at (-2,-3.6) {erros / resíduos\\do modelo 1};
\node[box, fill=alertbg, draw=alertframe] (E2) at (2,-3.6) {erros / resíduos\\do ensemble};

\node[box, fill=checkpointbg, draw=checkpointframe, minimum width=3.8cm] (F) at (0,-5.4)
{Soma ponderada\\dos modelos};

\node[box, fill=boxbg, draw=boxblue, minimum width=3.2cm] (Y) at (0,-7)
{previsão final};

\draw[arr] (D) -- (M1);
\draw[arr] (M1) -- (E1);
\draw[arr] (E1) -- (M2);
\draw[arr] (M2) -- (E2);
\draw[arr] (E2) -- (M3);

\draw[arr] (M1) -- (F);
\draw[arr] (M2) -- (F);
\draw[arr] (M3) -- (F);
\draw[arr] (F) -- (Y);

\node[font=\scriptsize, text=black!65, align=center] at (0,-8.1)
{Boosting: cada novo modelo é treinado para reduzir os erros deixados pelos anteriores.};
\end{tikzpicture}
\caption{Arquitetura conceitual do boosting. Diferentemente do bagging, os modelos não são treinados de forma independente. Eles são ajustados em sequência, de modo que cada novo modelo tenta corrigir erros ou resíduos deixados pelo ensemble acumulado.}
\label{fig:boosting}
\end{figure}

\begin{alertabox}{Boosting também pode sofrer overfitting}
Embora o boosting seja frequentemente apresentado como uma técnica de redução de viés, modelos boosted também podem sofrer overfitting se forem muito profundos, numerosos ou pouco regularizados. Implementações modernas como XGBoost e LightGBM incorporam regularização e amostragem de dados por esse motivo. Sempre avalie o modelo boosted tanto no treino quanto na validação.
\end{alertabox}

\subsection{Stacking: combinação aprendida de modelos}

O \emph{stacking} (\textit{stacked generalization}) combina modelos base de forma diferente. Em vez de usar uma regra fixa, como votação ou média, ele treina um novo modelo para aprender como combinar as previsões dos modelos base. Esse novo modelo é chamado de \emph{meta-modelo}.

A ideia é usar modelos base diversos, por exemplo uma árvore de decisão, uma regressão logística e um k-NN. Cada um produz uma previsão para o mesmo exemplo. Essas previsões passam a formar uma nova representação dos dados. O meta-modelo recebe essa representação como entrada e aprende a produzir a previsão final.

Em geral, o meta-modelo não deve ser treinado usando previsões feitas sobre os mesmos exemplos usados para ajustar os modelos base, pois isso pode causar vazamento de informação e overfitting. Uma prática comum é gerar previsões fora da amostra, por validação cruzada, para treinar o meta-modelo.

A Figura~\ref{fig:stacking} mostra a lógica do stacking. A diferença em relação ao bagging e ao boosting está na forma de combinação. No bagging, a combinação é uma votação ou média. No boosting, a combinação é acumulada sequencialmente. No stacking, a combinação é aprendida por outro modelo. Esse meta-modelo observa as previsões dos modelos base e aprende, por exemplo, quando confiar mais em um
classificador do que em outro. Por isso, stacking é especialmente útil quando os modelos base têm comportamentos complementares.

\begin{figure}[htb]
\centering
\begin{tikzpicture}[
    box/.style={draw, rounded corners=3pt, thick, minimum width=2.5cm,
                minimum height=0.8cm, align=center, font=\scriptsize},
    arr/.style={-{Stealth[length=5pt]}, thick},
    lbl/.style={font=\scriptsize, fill=white, inner sep=1pt}
]
\node[box, fill=statecolor, draw=agentcolor] (D) at (0,0) {Dados $\mathcal{D}$};

\node[box, fill=workedbg, draw=envcolor] (M1) at (-4,-1.5) {modelo base 1\\árvore};
\node[box, fill=workedbg, draw=envcolor] (M2) at (0,-1.5) {modelo base 2\\regressão};
\node[box, fill=workedbg, draw=envcolor] (M3) at (4,-1.5) {modelo base 3\\k-NN};

\node[box, fill=checkpointbg, draw=checkpointframe] (P1) at (-4,-3.3) {previsão 1};
\node[box, fill=checkpointbg, draw=checkpointframe] (P2) at (0,-3.3) {previsão 2};
\node[box, fill=checkpointbg, draw=checkpointframe] (P3) at (4,-3.3) {previsão 3};

\node[box, fill=alertbg, draw=alertframe, minimum width=4.2cm] (Z) at (0,-4.9)
{nova representação\\com previsões dos modelos base};

\node[box, fill=workedbg, draw=envcolor, minimum width=3.3cm] (Meta) at (0,-6.5)
{meta-modelo};

\node[box, fill=boxbg, draw=boxblue, minimum width=3.2cm] (Y) at (0,-8.1)
{previsão final};

\draw[arr] (D) -- (M1);
\draw[arr] (D) -- (M2);
\draw[arr] (D) -- (M3);

\draw[arr] (M1) -- (P1);
\draw[arr] (M2) -- (P2);
\draw[arr] (M3) -- (P3);

\draw[arr] (P1) -- (Z);
\draw[arr] (P2) -- (Z);
\draw[arr] (P3) -- (Z);

\draw[arr] (Z) -- (Meta);
\draw[arr] (Meta) -- (Y);

\node[font=\scriptsize, text=black!65, align=center] at (0,-9.7)
{Stacking: o modo de combinar os modelos base também é aprendido a partir dos dados.};
\end{tikzpicture}
\caption{Arquitetura conceitual do stacking. Modelos base possivelmente diferentes produzem previsões para os mesmos exemplos. Essas previsões são usadas como entrada para um meta-modelo, que aprende a combiná-las para produzir a previsão final.}
\label{fig:stacking}
\end{figure}

\subsection{Comparativo}

A Tabela~\ref{tab:ensembles} resume as principais diferenças entre os métodos de ensemble.

\begin{table}[htb]
\centering
\caption{Comparação entre métodos ensemble quanto ao componente do erro que atacam e à sua estratégia de combinação.}
\label{tab:ensembles}
\renewcommand{\arraystretch}{1.3}
\begin{tabular}{p{2.8cm}p{4.5cm}p{6.cm}}
\toprule
\textbf{Método} & \textbf{Reduz principalmente} & \textbf{Estratégia} \\
\midrule
Bagging & Variância & Modelos independentes treinados em amostras bootstrap; combinação por votação ou média \\
Random Forest & Variância & Bagging com árvores + aleatoriedade na seleção de atributos \\
AdaBoost & Viés & Modelos sequenciais; exemplos errados recebem maior peso \\
Gradient Boosting & Viés; variância controlada por regularização & Modelos sequenciais treinados para corrigir resíduos do ensemble acumulado \\
Stacking & Depende dos modelos base e do meta-modelo & Modelos base produzem previsões; um meta-modelo aprende como combiná-las \\
\bottomrule
\end{tabular}
\end{table}

\begin{checkpointbox}{Micro-checkpoint}
Árvores de decisão têm alta variância; modelos lineares simples têm alto viés. Como cada um dos métodos ensemble (bagging e boosting) ataca um desses dois problemas? Por que Random Forest funciona melhor que bagging de árvores simples?
\end{checkpointbox}

\begin{perguntabox}
Um conjunto de dados tem 1.000 exemplos e 20 atributos. Um único Decision Tree treinado até a folha (sem poda) obtém 100\% de acurácia no treino e 74\% na validação. Um Random Forest com 100 árvores obtém 88\% de acurácia na validação. \textbf{(a)} Por que o Random Forest generaliza melhor? \textbf{(b)} Qual mecanismo do Random Forest é responsável por diversificar as árvores e reduzir sua correlação? \textbf{(c)} Se você quisesse melhorar ainda mais o desempenho, como compararia Random Forest com Gradient Boosting para este problema?
\end{perguntabox}

\begin{perguntabox}
As notas de AR compararam Q-learning tabular (exato, mas não escala) com Q-learning aproximado (generaliza, mas sem garantia de $Q^*$). Trace um paralelo preciso com a comparação entre árvores de decisão (exatas sobre o treino, sem regularização) e Random Forests (generalizadas, mas sem garantia de ótimo global). O que esses dois casos têm em comum conceitualmente? O que os distingue?
\end{perguntabox}

%===========================================
\section{Redes Neurais Artificiais}
\label{sec:redes-neurais}

\subsection{Do perceptron à rede multicamadas}
\label{subsec:mlp}

O \emph{perceptron} é o bloco básico de construção das redes neurais. Ele recebe um vetor de entradas $\mathbf{x}$, computa uma combinação linear $z = \mathbf{w}^\top \mathbf{x} + b$, e aplica uma \emph{função de ativação} $\sigma$ para produzir uma saída:

\[
\text{saída} = \sigma(\mathbf{w}^\top \mathbf{x} + b)
\]

Com a função de ativação como degrau ($\sigma(z) = \mathbf{1}[z \geq 0]$), o perceptron é um classificador linear. Com a sigmoide como ativação, ele é equivalente à regressão logística. O perceptron simples é, portanto, um modelo linear, com a mesma limitação de não conseguir resolver problemas não linearmente separáveis.

O avanço crucial é o \emph{perceptron multicamadas} (\textit{multilayer perceptron}, MLP), que empilha múltiplas camadas de perceptrons. 
A Figura~\ref{fig:mlp} apresenta a estrutura de um perceptron multicamadas, ou MLP. Os valores de entrada $\mathbf{x}$ são propagados para uma camada oculta, onde cada neurônio combina todas as entradas por meio de pesos e vieses aprendíveis. Em seguida, aplica-se uma função de ativação não linear $\sigma$, produzindo a representação intermediária $\mathbf{h}$. Essa etapa é essencial: sem a não linearidade, a composição de camadas continuaria equivalente a uma única transformação linear. A camada de saída combina os valores ocultos e aplica uma função \textit{softmax}, produzindo uma distribuição de probabilidades sobre as classes. Assim, o MLP aprende simultaneamente uma representação dos dados e uma regra de decisão, permitindo modelar fronteiras mais complexas do que aquelas produzidas por modelos lineares.

\begin{figure}[ht]
\centering
\begin{tikzpicture}[
    neuron/.style={circle, draw=agentcolor, fill=statecolor, thick,
                   minimum size=0.9cm, font=\scriptsize, text=agentcolor},
    outneuron/.style={circle, draw=envcolor, fill=workedbg, thick,
                      minimum size=0.9cm, font=\scriptsize, text=envcolor},
    arr/.style={-{Stealth[length=4pt]}, thick, black!30},
    lbl/.style={font=\scriptsize}
]

% Camada de entrada (4 neurônios)
\foreach \i in {1,2,3,4} {
    \pgfmathsetmacro{\y}{(2.5-\i)*1.4}
    \node[neuron] (I\i) at (0, \y) {$x_\i$};
}

% Camada oculta (5 neurônios)
\foreach \i in {1,2,3,4,5} {
    \pgfmathsetmacro{\y}{(3-\i)*1.2}
    \node[neuron] (H\i) at (4, \y) {$h_\i$};
}

% Camada de saída (2 neurônios)
\foreach \i in {1,2} {
    \pgfmathsetmacro{\y}{(1.5-\i)*1.4}
    \node[outneuron] (O\i) at (8, \y) {$\hat{y}_\i$};
}

% Conexões entrada -> oculta
\foreach \i in {1,2,3,4} {
    \foreach \j in {1,2,3,4,5} {
        \draw[arr] (I\i) -- (H\j);
    }
}

% Conexões oculta -> saída
\foreach \i in {1,2,3,4,5} {
    \foreach \j in {1,2} {
        \draw[arr] (H\i) -- (O\j);
    }
}

% Labels
\node[lbl, text=agentcolor, below=0.4cm of I4] {camada de entrada};
\node[lbl, text=agentcolor, below=0.4cm of H5] {camada oculta};
\node[lbl, text=envcolor, below=0.4cm of O2]   {camada de saída};

% Notações matemáticas
\node[lbl, text=black!60, align=center] at (2, -4.3)
    {$\mathbf{h} = \sigma(W^{(1)}\mathbf{x} + \mathbf{b}^{(1)})$};
\node[lbl, text=black!60, align=center] at (6, -4.3)
    {$\hat{\mathbf{y}} = \text{softmax}(W^{(2)}\mathbf{h} + \mathbf{b}^{(2)})$};

\end{tikzpicture}
\caption{Arquitetura de um perceptron multicamadas (MLP) com 4 entradas, 5 neurônios ocultos e 2 saídas. Cada conexão possui um peso aprendível. A camada oculta aplica uma transformação não linear $\sigma$, permitindo ao modelo aprender fronteiras de decisão não lineares.}
\label{fig:mlp}
\end{figure}

O poder de representação do MLP vem da não-linearidade: sem funções de ativação não lineares, qualquer composição de camadas seria equivalente a uma única camada linear. Com não-linearidade, camadas ocultas podem aprender \emph{representações intermediárias} dos dados, uma forma automática de engenharia de características.

\begin{notabox}{Por que redes neurais são universais}
O \emph{teorema de aproximação universal} (\citep{hornik1991approximation,cybenko1989approximation}) afirma que um MLP com uma única camada oculta suficientemente larga pode aproximar qualquer função contínua com precisão arbitrária. Isso significa que a classe de hipóteses dos MLPs é extraordinariamente expressiva. A questão prática, no entanto, não é \emph{se} um MLP pode representar a função, mas se ele consegue aprendê-la a partir de dados finitos e com gradiente descendente.
\end{notabox}

\subsection{Funções de ativação}
\label{subsec:ativacao}

A escolha da função de ativação influencia fortemente a capacidade e a eficiência do treinamento. As mais usadas são:

\begin{itemize}
    \item Sigmoide: $\sigma(z) = \frac{1}{1+e^{-z}} \in (0,1)$. Historicamente importante, mas sofre de \textit{vanishing gradient} em redes profundas.
    \item Tanh: $\tanh(z) = \frac{e^z - e^{-z}}{e^z + e^{-z}} \in (-1,1)$. Similar à sigmoide, mas centrada em zero; geralmente preferível.
    \item ReLU (\textit{Rectified Linear Unit}): $\text{ReLU}(z) = \max(0, z)$. Simples, eficiente, e o padrão atual na maioria das arquiteturas. Não sofre de \textit{vanishing gradient} para $z > 0$.
    \item Softmax (camada de saída para classificação multiclasse): $\text{softmax}(\mathbf{z})_k = \frac{e^{z_k}}{\sum_j e^{z_j}}$. Normaliza as saídas para uma distribuição de probabilidade sobre $K$ classes.
\end{itemize}

\begin{alertabox}{O problema do gradiente que desaparece}
A sigmoide comprime qualquer entrada para $(0,1)$, e sua derivada máxima é $0{,}25$. Em uma rede profunda, o gradiente é multiplicado por essa derivada em cada camada durante a retropropagação. Com 10 camadas, o gradiente é multiplicado por no máximo $0{,}25^{10} \approx 10^{-6}$: torna-se praticamente zero. Isso faz com que os pesos das primeiras camadas aprendam muito lentamente ou nada. Esse problema, chamado de \emph{vanishing gradient}, foi a principal barreira ao treinamento de redes profundas antes da adoção da ReLU.
\end{alertabox}

\subsection{Retropropagação}
\label{subsec:backprop}

Como treinar os parâmetros de uma rede neural? O mesmo princípio do gradiente descendente se aplica, mas agora é necessário calcular o gradiente do custo em relação a \emph{todos} os pesos da rede. O algoritmo que faz isso de forma eficiente é a \emph{retropropagação} (\textit{backpropagation}).

A retropropagação não é um algoritmo de otimização; é um algoritmo de cálculo de gradientes. Ele aplica a \emph{regra da cadeia do cálculo} de forma sistemática, propagando o erro da camada de saída de volta para a camada de entrada:

\[
\frac{\partial J}{\partial w^{(l)}_{jk}} = \frac{\partial J}{\partial a^{(l)}_j} \cdot \frac{\partial a^{(l)}_j}{\partial w^{(l)}_{jk}}
\]

onde $a^{(l)}_j$ é a ativação do neurônio $j$ na camada $l$ e $w^{(l)}_{jk}$ é o peso que conecta o neurônio $k$ da camada $l-1$ ao neurônio $j$ da camada $l$.

A Figura~\ref{fig:backprop} resume as duas fases computacionais usadas no treinamento de uma rede neural simples. No passo \textit{forward}, os valores fluem da entrada $x$ para a representação intermediária $h$, depois para a predição $\hat{y}$ e, por fim, para a função de perda $J$. Essa etapa calcula o erro produzido pelo modelo para os parâmetros atuais. No passo \textit{backward}, o fluxo ocorre no sentido inverso: a partir da perda, calculam-se derivadas parciais sucessivas em relação às saídas intermediárias, como $\frac{\partial J}{\partial \hat{y}}$, $\frac{\partial J}{\partial h}$ e $\frac{\partial J}{\partial x}$. Essas derivadas são propagadas pela regra da cadeia e permitem calcular quanto cada peso e viés contribuiu para o erro. Assim, a retropropagação não é um novo modelo, mas um procedimento eficiente para obter os gradientes necessários à atualização dos parâmetros.

\begin{figure}[ht]
\centering
\begin{tikzpicture}[
    neuron/.style={circle, draw=agentcolor, fill=statecolor, thick,
                   minimum size=1.0cm, font=\small, text=agentcolor},
    arr/.style={-{Stealth[length=5pt]}, thick},
    barr/.style={-{Stealth[length=5pt]}, thick, dashed},
    lbl/.style={font=\scriptsize, fill=white, inner sep=1.5pt}
]

\node[neuron] (x) at (0, 0)   {$x$};
\node[neuron] (h) at (3.5, 0) {$h$};
\node[neuron] (y) at (7, 0)   {$\hat{y}$};

\node[circle, draw=alertframe, fill=alertbg, thick,
      minimum size=0.8cm, font=\small, text=alertframe] (L) at (10.5, 0) {$J$};

% Forward
\draw[arr, agentcolor] (x) -- (h) node[lbl, above, midway] {$w^{(1)}, b^{(1)}$};
\draw[arr, agentcolor] (h) -- (y) node[lbl, above, midway] {$w^{(2)}, b^{(2)}$};
\draw[arr, agentcolor] (y) -- (L) node[lbl, above, midway] {perda};

% Backward
\draw[barr, mlred] (L) to[bend left=25]
    node[lbl, below=2pt, midway, text=mlred]
    {$\frac{\partial J}{\partial \hat{y}}$}
(y);

\draw[barr, mlred] (y) to[bend left=25]
    node[lbl, below=2pt, midway, text=mlred]
    {$\frac{\partial J}{\partial h}$}
(h);

\draw[barr, mlred] (h) to[bend left=25]
    node[lbl, below=2pt, midway, text=mlred]
    {$\frac{\partial J}{\partial x}$}
(x);

\node[font=\scriptsize, text=agentcolor] at (3.5,  1.3) {\faArrowRight\ passo forward};
\node[font=\scriptsize, text=mlred]      at (3.5, -1.3) {\faArrowLeft\ passo backward};

\end{tikzpicture}
\caption{Intuição da retropropagação. No passo \textit{forward}, os dados são propagados da esquerda para a direita para computar a perda $J$. No passo \textit{backward}, os gradientes são propagados da direita para a esquerda via regra da cadeia, calculando $\frac{\partial J}{\partial w}$ para cada peso.}
\label{fig:backprop}
\end{figure}

A retropropagação calcula os gradientes de forma exata e eficiente em $O(W)$, onde $W$ é o número total de pesos da rede. Uma vez calculados, os gradientes são usados no gradiente descendente estocástico (SGD) para atualizar os pesos:

\[
w^{(l)}_{jk} \;\leftarrow\; w^{(l)}_{jk} - \alpha \cdot \frac{\partial J}{\partial w^{(l)}_{jk}}
\]

\begin{notabox}{Retropropagação não é magia}
Um equívoco comum é imaginar que a retropropagação ``ajusta os pesos inteligentemente''. Na realidade, ela faz uma coisa simples: aplicar a regra da cadeia sistematicamente para calcular derivadas parciais. O que há de elegante é a eficiência: em vez de estimar gradientes por diferenças finitas (custoso), a retropropagação os calcula de forma exata em tempo linear no número de pesos.

A atualização dos pesos em si é o mesmo gradiente descendente da regressão linear. A única novidade é o cálculo eficiente dos gradientes numa função composta por muitas camadas.
\end{notabox}

\subsection{Aprendizado profundo: uma visão geral}
\label{subsec:deep-learning}

O \emph{aprendizado profundo} (\textit{deep learning}) é o nome dado ao uso de redes neurais com muitas camadas ocultas. A profundidade permite que as camadas intermediárias aprendam representações hierárquicas dos dados: na visão computacional, por exemplo, as primeiras camadas detectam bordas, as intermediárias detectam formas, e as últimas detectam objetos.

O sucesso do aprendizado profundo a partir de 2012 (com o AlexNet em ImageNet) foi impulsionado por três fatores:

\begin{itemize}
    \item Dados: disponibilidade de conjuntos de dados maciços (ImageNet, Common Crawl, etc.).
    \item Computação: GPUs que permitem treinamento de redes grandes em tempo razoável.
    \item Algoritmos: melhores funções de ativação (ReLU), regularização (Dropout), normalização de lotes (\textit{batch normalization}), e arquiteturas especializadas.
\end{itemize}

As arquiteturas mais influentes incluem:

\begin{itemize}
    \item Redes convolucionais (CNNs): especializadas em dados com estrutura espacial (imagens).
    \item Redes recorrentes (RNNs / LSTMs): especializadas em sequências (texto, áudio, séries temporais).
    \item Transformers: arquitetura baseada em atenção, dominante em processamento de linguagem natural; base dos modelos GPT, BERT e Claude.
\end{itemize}

\begin{notabox}{Aprendizado profundo e engenharia de características}
Uma das contribuições mais importantes do aprendizado profundo é automatizar a engenharia de características: em vez de extrair manualmente $f_i(x)$, a rede aprende automaticamente quais características são relevantes como parte do treinamento. Isso explica por que modelos profundos, com dados suficientes, frequentemente superam modelos ``rasos'' com características manuais cuidadosamente projetadas.
\end{notabox}

\begin{alertabox}{Redes profundas não são solução universal}
Redes neurais profundas têm alta capacidade, mas essa capacidade só é útil quando há dados, computação e regularização suficientes. Em bases pequenas ou tabulares, modelos mais simples (como regressão logística, Random Forests ou Gradient Boosting) frequentemente apresentam desempenho equivalente ou superior, com muito menos custo computacional e maior interpretabilidade. A escolha do modelo deve ser guiada pelo problema, não pela complexidade.
\end{alertabox}

\begin{checkpointbox}{Micro-checkpoint}
Explique em uma frase o que a retropropagação faz. Descreva o papel da função de ativação numa rede neural: o que aconteceria se todas as ativações fossem lineares? Qual arquitetura de rede neural seria mais adequada para classificar imagens de raio-X médico?
\end{checkpointbox}

\begin{perguntabox}
Considere uma rede neural com uma camada oculta de 10 neurônios ReLU treinada para classificação binária (spam/não-spam). O modelo obtém 99\% de acurácia no treino mas apenas 65\% na validação. \textbf{(a)} Qual problema está ocorrendo? \textbf{(b)} Cite três técnicas específicas de treinamento de redes neurais que poderiam melhorar a generalização. \textbf{(c)} Como o conceito de regularização (Seção~\ref{sec:generalizacao}) se aplica a redes neurais?
\end{perguntabox}

%===========================================
\section{Aprendizado Não Supervisionado}
\label{sec:nao-supervisionado}

\subsection{Quando não temos rótulos}
\label{subsec:nsup-intro}

Nos problemas de aprendizado supervisionado, cada exemplo de treinamento vem acompanhado de um rótulo $y$. Na prática, rotular dados é caro, demorado e muitas vezes impraticável em larga escala. O \emph{aprendizado não supervisionado} aborda o problema de extrair estrutura de dados \emph{sem rótulos}.

As tarefas mais comuns são:

\begin{itemize}
    \item Agrupamento (\textit{clustering}): encontrar grupos naturais nos dados.
    \item Redução de dimensionalidade: representar os dados em um espaço de menor dimensão, preservando estrutura relevante.
    \item Estimação de densidade: modelar a distribuição $p(\mathbf{x})$.
\end{itemize}

\subsection{\texorpdfstring{$k$}{k}-means}
\label{subsec:kmeans}

O algoritmo $k$-means é o método de agrupamento mais simples e amplamente usado. Dado um conjunto de pontos $\{\mathbf{x}^{(1)}, \ldots, \mathbf{x}^{(n)}\}$ e um número de grupos $k$, o $k$-means encontra $k$ \emph{centróides} $\{\boldsymbol{\mu}_1, \ldots, \boldsymbol{\mu}_k\}$ que minimizam a soma das distâncias quadráticas de cada ponto ao seu centróide mais próximo:

\[
J(\{C_i\}, \{\boldsymbol{\mu}_i\}) = \sum_{i=1}^k \sum_{\mathbf{x} \in C_i} \|\mathbf{x} - \boldsymbol{\mu}_i\|^2
\]

O algoritmo alterna entre dois passos até convergência:

\begin{algorithm}[htb]
\caption{$k$-means}
\begin{algorithmic}[1]
\STATE Inicializar $k$ centróides $\boldsymbol{\mu}_1, \ldots, \boldsymbol{\mu}_k$ (e.g., $k$ pontos aleatórios do conjunto)
\REPEAT
    \STATE \textbf{Passo E (atribuição)}: atribuir cada ponto $\mathbf{x}^{(i)}$ ao grupo cujo centróide é mais próximo:
    \[
    c^{(i)} \leftarrow \arg\min_j \|\mathbf{x}^{(i)} - \boldsymbol{\mu}_j\|^2
    \]
    \STATE \textbf{Passo M (atualização)}: recomputar cada centróide como a média dos pontos atribuídos a ele:
    \[
    \boldsymbol{\mu}_j \leftarrow \frac{1}{|C_j|} \sum_{\mathbf{x} \in C_j} \mathbf{x}
    \]
\UNTIL{as atribuições $c^{(i)}$ não mudarem}
\RETURN centróides $\{\boldsymbol{\mu}_j\}$ e atribuições $\{c^{(i)}\}$
\end{algorithmic}
\label{alg:kmeans}
\end{algorithm}

A Figura~\ref{fig:kmeans} ilustra o funcionamento do algoritmo $k$-means para $k=3$. No painel inicial, os pontos ainda não possuem grupos definidos, e os centróides começam em posições escolhidas aleatoriamente, indicadas pelas estrelas. A cada iteração, o algoritmo executa dois passos: primeiro, atribui cada ponto ao centróide mais próximo; depois, reposiciona cada centróide na média dos pontos atribuídos ao seu grupo. Esse processo se repete até que as atribuições deixem de mudar de forma relevante. No painel final, os pontos aparecem separados em três agrupamentos estáveis, cada um representado por uma cor e por seu centróide final. A figura também sugere uma limitação importante: como o resultado depende dos centróides iniciais, diferentes inicializações podem levar a agrupamentos diferentes.

\begin{figure}[ht]
\centering
\begin{tikzpicture}[
    pt/.style={circle, minimum size=5pt, inner sep=0},
    cent/.style={star, star points=5, draw=black, thick, minimum size=10pt, inner sep=0}
]

% Iteração 0 (pontos sem agrupamento)
\begin{scope}[xshift=0cm]
\node[font=\small\bfseries, text=black!70] at (2, 4.5) {Inicial};
% pontos (cluster 1)
\foreach \xy in {(0.5,3),(1.0,3.5),(0.8,2.5),(1.3,3.2),(0.4,2.8)} {
    \node[pt, fill=black!30] at \xy {}; }
% pontos (cluster 2)
\foreach \xy in {(3.0,3.3),(3.5,2.8),(2.8,2.5),(3.2,3.8),(3.8,3.1)} {
    \node[pt, fill=black!30] at \xy {}; }
% pontos (cluster 3)
\foreach \xy in {(1.8,0.8),(2.3,1.2),(1.5,1.5),(2.6,0.6),(2.0,1.8)} {
    \node[pt, fill=black!30] at \xy {}; }
% centróides iniciais (aleatórios)
\node[cent, fill=mlred!70]   at (1.5, 1.0) {};
\node[cent, fill=agentcolor!70] at (2.5, 3.5) {};
\node[cent, fill=envcolor!70]  at (0.8, 2.0) {};
\end{scope}

% Seta
\draw[-{Stealth[length=6pt]}, thick] (4.5, 2.5) -- (5.5, 2.5)
    node[midway, above, font=\tiny] {iter. 1--3};

% Iteração final (pontos agrupados)
\begin{scope}[xshift=6cm]
\node[font=\small\bfseries, text=black!70] at (2, 4.5) {Convergido};
\foreach \xy in {(0.5,3),(1.0,3.5),(0.8,2.5),(1.3,3.2),(0.4,2.8)} {
    \node[pt, fill=agentcolor!70] at \xy {}; }
\foreach \xy in {(3.0,3.3),(3.5,2.8),(2.8,2.5),(3.2,3.8),(3.8,3.1)} {
    \node[pt, fill=mlred!70] at \xy {}; }
\foreach \xy in {(1.8,0.8),(2.3,1.2),(1.5,1.5),(2.6,0.6),(2.0,1.8)} {
    \node[pt, fill=envcolor!70] at \xy {}; }
% centróides finais
\node[cent, fill=agentcolor, draw=black] at (0.8, 3.1) {};
\node[cent, fill=mlred,   draw=black]   at (3.3, 3.1) {};
\node[cent, fill=envcolor, draw=black]  at (2.0, 1.2) {};
\end{scope}
\end{tikzpicture}
\caption{O algoritmo $k$-means com $k=3$. Partindo de centróides iniciais aleatórios (estrelas), o algoritmo alterna entre atribuição de pontos ao centróide mais próximo e recálculo dos centróides, convergindo para um agrupamento estável.}
\label{fig:kmeans}
\end{figure}

\subsubsection{Exemplo resolvido}

Aplique $k$-means com $k=2$ ao conjunto $\{(1,1), (1,2), (5,4), (5,5)\}$. Centróides iniciais: $\boldsymbol{\mu}_1 = (1,1)$ e $\boldsymbol{\mu}_2 = (5,4)$.

\paragraph{Iteração 1 -- Passo E:}
\begin{itemize}[nosep]
    \item $(1,1)$: dist. a $\boldsymbol{\mu}_1 = 0$; dist. a $\boldsymbol{\mu}_2 = \sqrt{32} \approx 5{,}7$. Atribui ao grupo 1.
    \item $(1,2)$: dist. a $\boldsymbol{\mu}_1 = 1$; dist. a $\boldsymbol{\mu}_2 = \sqrt{25} = 5$. Atribui ao grupo 1.
    \item $(5,4)$: dist. a $\boldsymbol{\mu}_1 = \sqrt{25} = 5$; dist. a $\boldsymbol{\mu}_2 = 0$. Atribui ao grupo 2.
    \item $(5,5)$: dist. a $\boldsymbol{\mu}_1 = \sqrt{32} \approx 5{,}7$; dist. a $\boldsymbol{\mu}_2 = 1$. Atribui ao grupo 2.
\end{itemize}

\paragraph{Iteração 1 -- Passo M:}
\[
\boldsymbol{\mu}_1 = \frac{(1,1) + (1,2)}{2} = (1,\; 1{,}5) \qquad \boldsymbol{\mu}_2 = \frac{(5,4) + (5,5)}{2} = (5,\; 4{,}5)
\]

\paragraph{Iteração 2 -- Passo E:} as atribuições não mudam. Algoritmo converge.

\paragraph{Interpretação:} os dois grupos naturais $\{(1,1),(1,2)\}$ e $\{(5,4),(5,5)\}$ foram corretamente identificados em duas iterações.

\begin{alertabox}{Limitações do $k$-means}
O $k$-means tem limitações importantes: \textbf{(a)} requer que $k$ seja fornecido a priori; \textbf{(b)} converge para um mínimo local que depende da inicialização (o algoritmo \textit{k-means++} melhora a inicialização); \textbf{(c)} assume grupos de forma esférica e variância similar; falha com grupos alongados, em forma de lua crescente, ou com densidade muito diferente.
\end{alertabox}

\subsection{Análise de Componentes Principais (PCA)}
\label{subsec:pca}

A \emph{análise de componentes principais} (PCA, \textit{Principal Component Analysis}) é o método de redução de dimensionalidade mais clássico. Dado um conjunto de dados em $\mathbb{R}^d$, o PCA encontra uma projeção linear para um espaço de dimensão $k < d$ que \emph{maximiza a variância preservada} (equivalentemente, minimiza o erro de reconstrução).

A Figura~\ref{fig:pca} mostra a intuição geométrica da Análise de Componentes Principais em duas dimensões. A nuvem de pontos possui uma direção dominante de espalhamento: os dados variam mais ao longo do eixo indicado por PC$_1$. Essa é a primeira componente principal. A segunda componente, PC$_2$, é perpendicular à primeira e captura a maior parte da variação que ainda resta depois de PC$_1$. Quando projetamos os dados apenas sobre PC$_1$, substituímos cada ponto bidimensional por uma única coordenada ao longo dessa direção. Assim, reduzimos a dimensionalidade de 2 para 1 preservando o máximo possível da variância dos dados. A figura também ajuda a lembrar que PCA não usa rótulos de classe: ele procura direções de maior variância, não necessariamente direções de melhor separação entre grupos.

\begin{figure}[ht]
\centering
\begin{tikzpicture}
\begin{axis}[
    width=9cm, height=7cm,
    xlabel={$x_1$}, ylabel={$x_2$},
    xlabel style={font=\small},
    ylabel style={font=\small},
    xmin=-4, xmax=4, ymin=-3, ymax=3,
    xtick=\empty, ytick=\empty,
    clip=false
]
% Pontos em uma elipse inclinada
\foreach \a in {0,15,...,345} {
    \addplot[only marks, mark=*, mark size=1.5pt, agentcolor]
    coordinates {({2.5*cos(\a) - 0.8*sin(\a)}, {0.6*cos(\a) + 0.3*sin(\a)})};
}

% PC1 (direção de maior variância)
\addplot[thick, mlred, ->, domain=-1:1]
    ({2.8*x}, {1.1*x})
    node[pos=1, right, font=\small, text=mlred] {PC$_1$};

% PC2 (direção perpendicular)
\addplot[thick, envcolor, ->, domain=-0.5:0.5]
    ({-0.9*x}, {2.3*x})
    node[pos=1, above, font=\small, text=envcolor] {PC$_2$};

\node[font=\scriptsize, text=black!60, align=center] at (axis cs: 0, -2.5)
{PC$_1$ captura a maior variância; PC$_2$ o restante ortogonal};
\end{axis}
\end{tikzpicture}
\caption{PCA em 2D. Os dados formam uma nuvem elíptica. A primeira componente principal (PC$_1$) aponta na direção de maior variância; a segunda (PC$_2$) é perpendicular a ela. Projetar os dados sobre apenas PC$_1$ reduz a dimensão de 2 para 1, com mínima perda de informação.}
\label{fig:pca}
\end{figure}

Intuitivamente, o PCA encontra as direções do espaço de características ao longo das quais os dados variam mais. Projetar os dados sobre as $k$ primeiras componentes principais retém a maior parte da variância com apenas $k$ dimensões.

O PCA tem múltiplas aplicações práticas:

\begin{itemize}
    \item Visualização: projetar dados de alta dimensão para 2D ou 3D.
    \item Pré-processamento: reduzir dimensionalidade antes de aplicar outro algoritmo de AM, acelerando o treinamento.
    \item Remoção de ruído: as componentes de menor variância frequentemente capturam ruído; descartá-las pode melhorar o sinal.
    \item Compressão: representar imagens ou documentos de forma compacta.
\end{itemize}

\begin{notabox}{PCA e redes neurais}
O PCA computa uma projeção linear ótima. Autocodificadores (\textit{autoencoders}) são redes neurais treinadas para comprimir e reconstruir dados. Com ativações lineares, um autocodificador de uma camada é equivalente ao PCA. Com ativações não lineares, o autocodificador aprende projeções não lineares, uma generalização poderosa do PCA.
\end{notabox}

\begin{checkpointbox}{Micro-checkpoint}
Qual a diferença fundamental entre aprendizado supervisionado e não supervisionado? O que o $k$-means minimiza? Por que é necessário especificar $k$ a priori, e como você poderia escolher um bom valor de $k$ sem conhecimento prévio?
\end{checkpointbox}

\begin{perguntabox}
Uma empresa de varejo quer segmentar seus clientes em grupos para personalizar campanhas de marketing. Os dados disponíveis incluem: valor médio de compra, frequência de visitas mensais, categorias de produtos comprados (como um vetor de 50 dimensões representando proporção de compras em cada categoria) e informações demográficas. \textbf{(a)} Argumente por que esta é uma tarefa de aprendizado não supervisionado, e não de aprendizado supervisionado. \textbf{(b)} Quais são as limitações do $k$-means para este problema? \textbf{(c)} Por que seria útil aplicar PCA antes do $k$-means neste caso?
\end{perguntabox}

%===========================================
\section{Resumo}
\label{sec:resumo}

O AM é o estudo de algoritmos que aprendem a partir de dados. O framework supervisionado formaliza o aprendizado como minimização de um risco empírico num espaço de hipóteses, sujeito à restrição implícita de generalização. O dilema viés-variância capta a tensão fundamental entre modelos simples (que generalizam mas erram sistematicamente) e modelos complexos (que ajustam bem o treino mas podem não generalizar). A regularização, a validação cruzada e a divisão adequada dos dados são os instrumentos práticos para navegar esse dilema.

Modelos lineares são simples, eficientes e interpretáveis, mas limitados. Árvores de decisão são interpretáveis mas instáveis; métodos ensemble as tornam robustas. Redes neurais profundas são universalmente expressivas e aprendem características automaticamente, ao custo de requisitos maiores de dados, computação e conhecimento de treinamento. Cada família de modelos tem seu lugar, e a escolha deve ser guiada pela natureza do problema, pelo volume e qualidade dos dados, e pelos requisitos de interpretabilidade e desempenho.

A lista a seguir sintetiza os algoritmos de AM estudados nesta leitura, destacando paradigma, tipo de função aprendida, perfil viés-variância e observações práticas.

\begin{itemize}
    \item \textbf{Regressão linear}
    \begin{itemize}
        \item \textbf{Paradigma:} supervisionado.
        \item \textbf{O que aprende:} uma função linear de predição.
        \item \textbf{Viés/variância:} alto viés e baixa variância.
        \item \textbf{Escalabilidade/notas:} admite solução analítica quando $d$ é pequeno; pode ser treinada com SGD quando $d$ é grande.
    \end{itemize}

    \item \textbf{Regressão logística}
    \begin{itemize}
        \item \textbf{Paradigma:} supervisionado.
        \item \textbf{O que aprende:} uma fronteira linear de decisão para classificação.
        \item \textbf{Viés/variância:} alto viés e baixa variância.
        \item \textbf{Escalabilidade/notas:} eficiente, interpretável e base para o classificador softmax.
    \end{itemize}

    \item \textbf{Árvore de decisão}
    \begin{itemize}
        \item \textbf{Paradigma:} supervisionado.
        \item \textbf{O que aprende:} um particionamento hierárquico do espaço de entrada.
        \item \textbf{Viés/variância:} baixo viés e alta variância.
        \item \textbf{Escalabilidade/notas:} interpretável, mas instável; tende a overfitting sem poda ou outras formas de regularização.
    \end{itemize}

    \item \textbf{Random Forest}
    \begin{itemize}
        \item \textbf{Paradigma:} supervisionado.
        \item \textbf{O que aprende:} um ensemble de árvores de decisão.
        \item \textbf{Viés/variância:} reduz a variância em relação a uma única árvore.
        \item \textbf{Escalabilidade/notas:} robusto, pouco sensível a hiperparâmetros e escalável por paralelismo entre árvores.
    \end{itemize}

    \item \textbf{Gradient Boosting}
    \begin{itemize}
        \item \textbf{Paradigma:} supervisionado.
        \item \textbf{O que aprende:} um ensemble sequencial de modelos fracos.
        \item \textbf{Viés/variância:} pode reduzir viés e variância.
        \item \textbf{Escalabilidade/notas:} muito poderoso, mas mais sensível a hiperparâmetros.
    \end{itemize}

    \item \textbf{MLP / Deep Learning}
    \begin{itemize}
        \item \textbf{Paradigma:} supervisionado.
        \item \textbf{O que aprende:} representações hierárquicas dos dados.
        \item \textbf{Viés/variância:} depende da arquitetura, do tamanho do modelo, da regularização e dos dados.
        \item \textbf{Escalabilidade/notas:} alta capacidade expressiva; requer muitos dados e treinamento computacionalmente custoso.
    \end{itemize}

    \item \textbf{$k$-means}
    \begin{itemize}
        \item \textbf{Paradigma:} não supervisionado.
        \item \textbf{O que aprende:} centróides que representam grupos nos dados.
        \item \textbf{Viés/variância:} não se aplica diretamente no mesmo sentido usado para modelos supervisionados.
        \item \textbf{Escalabilidade/notas:} sensível à inicialização e à escolha de $k$; assume implicitamente clusters aproximadamente esféricos.
    \end{itemize}

    \item \textbf{PCA}
    \begin{itemize}
        \item \textbf{Paradigma:} não supervisionado.
        \item \textbf{O que aprende:} uma projeção linear nas direções de maior variância.
        \item \textbf{Viés/variância:} não se aplica diretamente no mesmo sentido usado para modelos supervisionados.
        \item \textbf{Escalabilidade/notas:} fornece uma projeção linear ótima em termos de variância preservada; é eficiente e serve de base para muitas técnicas.
    \end{itemize}
\end{itemize}

%===========================================
\section{Atividades Intra-Classe}

As atividades abaixo devem ser resolvidas em grupo, com foco em explicar o raciocínio passo a passo, e não apenas em obter a resposta final.

\subsection{Aquecimento}

Antes das atividades principais, responda rapidamente:

\begin{enumerate}[label=(\alph*)]
    \item Qual a diferença entre um parâmetro (como os pesos $\mathbf{w}$) e um hiperparâmetro (como $\lambda$ na regularização)? O que é ajustado durante o treinamento e o que é ajustado durante a validação?
    \item Um modelo obteve erro de treino $= 0{,}02$ e erro de validação $= 0{,}85$. Qual problema você diagnosticaria? E se o erro de treino fosse $0{,}80$ e o de validação $0{,}82$?
    \item Cite a diferença entre MSE e entropia cruzada. Para qual tarefa cada uma é mais adequada?
    \item Explique a diferença entre bagging e boosting em uma frase cada.
    \item Por que inicializar todos os pesos de uma rede neural com zero é problemático? (\textit{Dica}: pense no que acontece com o gradiente se todos os neurônios de uma camada tiverem os mesmos pesos.)
    \item Em $k$-means, o que o algoritmo está minimizando? O que garante convergência?
\end{enumerate}

\begin{checkpointbox}{Explique como se fosse um professor}
Explique para um colega a diferença entre viés e variância usando apenas:
\begin{itemize}[nosep]
    \item a analogia de um arqueiro num campo de tiro;
    \item o comportamento do modelo com dados de treino vs. dados de validação;
    \item um exemplo concreto com regressão linear de grau 1 vs. grau 15;
    \item por que um modelo com alto viés e um com alta variância cometem erros muito diferentes.
\end{itemize}
Evite fórmulas. Se você não consegue explicar sem equações, revise a Seção~\ref{sec:generalizacao}.
\end{checkpointbox}

\subsection{Regressão linear e gradiente descendente}

Considere um problema de regressão com os dados a seguir:

\medskip
\begin{center}
\begin{tabular}{cccc}
\toprule
Exemplo & Área ($x_1$, m²) & Quartos ($x_2$) & Preço ($y$, R\$ mil) \\
\midrule
1 & 50 & 1 & 200 \\
2 & 80 & 2 & 300 \\
3 & 120& 3 & 450 \\
4 & 60 & 2 & 250 \\
\bottomrule
\end{tabular}
\end{center}

\begin{enumerate}[label=(\alph*)]
    \item Escreva a hipótese do modelo: $h(\mathbf{x}) = w_0 + w_1 x_1 + w_2 x_2$. Calcule o MSE para os pesos iniciais $w_0 = 0$, $w_1 = 3$, $w_2 = 10$.
    \item Calcule o gradiente $\frac{\partial J}{\partial w_1}$ e $\frac{\partial J}{\partial w_2}$ com os pesos acima. (Use $\frac{\partial J}{\partial w_j} = \frac{2}{n} \sum_i (\hat{y}_i - y_i) x_{ij}$.)
    \item Execute um passo de gradiente descendente com $\alpha = 0{,}0001$ e reporte os pesos atualizados.
    \item Após o passo, o MSE aumentou ou diminuiu? Por que isso era esperado?
\end{enumerate}

\subsection{Entropia e construção de árvores}

Considere um conjunto de 12 exemplos de e-mails: 8 spam e 4 legítimos. Dois atributos candidatos para a divisão raiz:

\begin{itemize}
    \item \textbf{Atributo A} (``oferta'' presente): $S_{\text{sim}} = \{6\text{spam}, 1\text{leg}\}$, $S_{\text{não}} = \{2\text{spam}, 3\text{leg}\}$.
    \item \textbf{Atributo B} (enviado de noite): $S_{\text{sim}} = \{4\text{spam}, 2\text{leg}\}$, $S_{\text{não}} = \{4\text{spam}, 2\text{leg}\}$.
\end{itemize}

\begin{enumerate}[label=(\alph*)]
    \item Calcule a entropia do conjunto raiz $H(S)$.
    \item Calcule o ganho de informação $\text{IG}(S, A)$ e $\text{IG}(S, B)$.
    \item Qual atributo o algoritmo de árvore escolheria? Justifique intuitivamente.
    \item O Atributo B tem IG igual a zero. O que isso diz sobre o poder discriminativo desse atributo? (Calcule para confirmar.)
\end{enumerate}

\subsection{Diagnóstico de generalização}

Para cada cenário abaixo, diagnostique o problema (underfitting, overfitting, ou adequado), identifique o provável desequilíbrio viés-variância e proponha uma ação corretiva específica.

\begin{enumerate}[label=(\alph*)]
    \item Um modelo de regressão linear simples ($d=1$, sem termos polinomiais) para prever o consumo de energia em função da temperatura, com erro de treino $= 320$ e erro de validação $= 335$.
    \item Uma Random Forest com 500 árvores, profundidade máxima ilimitada, treinada em 200 exemplos, com erro de treino $= 0$ e erro de validação $= 0{,}45$ (acurácia de 55\%).
    \item Uma regressão logística com regularização $\ell_2$, $\lambda = 1000$, com acurácia de treino $= 53\%$ e de validação $= 52\%$.
    \item Um MLP com 5 camadas ocultas de 1.000 neurônios, treinado por 500 épocas, com perda de treino $= 0{,}001$ e perda de validação $= 2{,}8$.
\end{enumerate}

\subsection{\texorpdfstring{$k$}{k}-means: execução manual}

Aplique o algoritmo $k$-means com $k=3$ ao conjunto de pontos em 2D (coordenadas inteiras):

\medskip
\begin{center}
\begin{tabular}{cc}
\toprule
Ponto & Coordenadas \\
\midrule
A & (1, 1) \\ B & (1, 2) \\ C & (2, 1) \\
D & (8, 8) \\ E & (9, 8) \\ F & (8, 9) \\
G & (4, 5) \\ H & (5, 4) \\ I & (5, 5) \\
\bottomrule
\end{tabular}
\end{center}

Centróides iniciais: $\boldsymbol{\mu}_1 = (1,1)$, $\boldsymbol{\mu}_2 = (4,5)$, $\boldsymbol{\mu}_3 = (8,8)$.

\begin{enumerate}[label=(\alph*)]
    \item Execute a iteração 1 completa: passo E (atribuição) e passo M (atualização dos centróides). Use distância euclidiana.
    \item Execute a iteração 2 e verifique se o algoritmo convergiu.
    \item O custo $J$ diminuiu da iteração 1 para a iteração 2? Por que o custo nunca aumenta entre iterações no $k$-means?
\end{enumerate}

\subsection{Perguntas para discussão em sala}

\begin{perguntabox}
Considere o seguinte cenário: você tem 100.000 e-mails rotulados (spam/legítimo) e precisa escolher entre três modelos: regressão logística, Random Forest e MLP com 3 camadas. Descreva o processo completo de desenvolvimento e avaliação, desde a divisão dos dados até a implantação do modelo. Que critérios usaria para a escolha final entre os três modelos?
\end{perguntabox}

\begin{perguntabox}
O dilema viés-variância implica que não existe um modelo universalmente melhor. No entanto, nas competições Kaggle de AM, Gradient Boosting (XGBoost/LightGBM) e redes neurais dominam consistentemente. Como reconciliar o teorema ``no free lunch'' com o sucesso prático desses métodos específicos em tantos problemas? (\textit{Dica}: considere que competições Kaggle não amostram uniformemente todos os problemas possíveis de AM; elas concentram certos tipos de dados, métricas e tarefas.)
\end{perguntabox}

\begin{perguntabox}
Um sistema de concessão de crédito usa um modelo de AM para decidir se um empréstimo deve ser aprovado. O modelo foi treinado com dados históricos de uma instituição financeira. \textbf{(a)} Quais riscos surgem se a hipótese de aprendizagem não for satisfeita nesse contexto? \textbf{(b)} Por que a acurácia pode ser uma métrica inadequada? \textbf{(c)} Que considerações éticas surgem quando os dados históricos refletem vieses sociais existentes?
\end{perguntabox}

\begin{perguntabox}
As notas de AR descrevem o Q-learning com aproximação linear como um gradiente descendente sobre o erro temporal. À luz do que foi estudado sobre AM, identifique: \textbf{(a)} o ``conjunto de treinamento'' equivalente no Q-learning; \textbf{(b)} o ``rótulo'' ($y$) equivalente; \textbf{(c)} de que forma o problema de generalização aparece no Q-learning com aproximação, e por que ele assume uma forma diferente do overfitting supervisionado clássico; \textbf{(d)} que tipo de regularização seria natural aplicar ao Q-learning com aproximação.
\end{perguntabox}

% ============================================================
\bibliographystyle{plainnat}
\bibliography{refs}
% ============================================================

\end{document}